\documentclass[12pt,a4paper]{article}

\usepackage[utf8]{inputenc}
\usepackage[T1]{fontenc}
\usepackage[french]{babel}
\addto\captionsfrench{%
  \renewcommand{\tablename}{Tableau}%
}
\usepackage{lmodern}
\usepackage{geometry}
\geometry{
  top=2.3cm,
  bottom=2.4cm,
  left=2.6cm,
  right=2.6cm
}

\usepackage{graphicx}
\graphicspath{{./}{images/}}
\usepackage{booktabs}
\usepackage{array}
\usepackage{tabularx}
\usepackage{caption}
\usepackage{subcaption}
\usepackage{float}
\usepackage{amsmath}
\usepackage{amssymb}
\usepackage{xcolor}
\usepackage{hyperref}
\usepackage{xurl}
\usepackage{fancyhdr}
\usepackage{titlesec}
\usepackage{setspace}
\usepackage{enumitem}
\usepackage{microtype}
\usepackage{pgfgantt}
\usepackage{pdflscape}
\usepackage[section]{placeins}

\onehalfspacing
\setlength{\parindent}{0.75cm}
\setlength{\parskip}{0.10cm}
\setlength{\headheight}{14pt}
\setlength{\textfloatsep}{10pt plus 2pt minus 2pt}
\setlength{\floatsep}{8pt plus 2pt minus 2pt}
\setlength{\intextsep}{8pt plus 2pt minus 2pt}
\setcounter{tocdepth}{2}
\setcounter{secnumdepth}{3}
\raggedbottom

\pagestyle{fancy}
\fancyhf{}
\lhead{\small TOUHAF Aya -- LISIC / ULCO}
\rhead{\small TER -- Master 1 Informatique}
\cfoot{\thepage}
\renewcommand{\headrulewidth}{0.4pt}

\titleformat{\section}
  {\large\bfseries}
  {\thesection.}
  {0.6em}
  {}
  [\titlerule]

\titleformat{\subsection}
  {\normalsize\bfseries}
  {\thesubsection.}
  {0.5em}
  {}

\titleformat{\subsubsection}
  {\normalsize\itshape}
  {\thesubsubsection.}
  {0.5em}
  {}

\hypersetup{
  colorlinks=true,
  linkcolor=blue!45!black,
  citecolor=green!35!black,
  urlcolor=blue!60!black,
  pdftitle={Vers une Intelligence Artificielle Explicable pour la Surveillance de l'État Physiologique de l'Athlète de Haut Niveau},
  pdfauthor={TOUHAF Aya}
}

\newcommand{\safeincludegraphics}[2][]{%
  \IfFileExists{#2}{%
    \includegraphics[#1]{#2}%
  }{%
    \fbox{%
      \begin{minipage}[c][4.2cm][c]{0.88\linewidth}
        \centering
        \textit{Figure non trouvée lors de la compilation}\\[0.2cm]
        \texttt{\detokenize{#2}}
      \end{minipage}
    }%
  }%
}

\newcommand{\HR}{\mathrm{HR}}
\newcommand{\BR}{\mathrm{BR}}
\newcommand{\RR}{\mathrm{RR}}

\newcommand{\openaccessref}[2]{%
  \nobreak\hspace{0.08em}%
  \href{#2}{\textcolor{green!35!black}{[#1]}}%
}


\begin{document}
\hypersetup{pageanchor=false}

% Page de titre

\begin{titlepage}
\thispagestyle{empty}
\centering

\vspace*{0.1cm}

\IfFileExists{images/logo_ulco.png}{%
  \includegraphics[width=0.20\textwidth]{images/logo_ulco.png}%
}{%
  {\Large\bfseries ULCO}%
}

\vspace{0.45cm}

{\Large\textsc{Université du Littoral Côte d'Opale}}\\[0.18cm]
{\large Laboratoire d'Informatique Signal et Image de la Côte d'Opale}\\[0.10cm]
{\large LISIC}\\[0.25cm]
{\normalsize Master Informatique -- Première année}

\vspace{0.75cm}

\rule{\linewidth}{0.8pt}\\[0.30cm]

{\LARGE\bfseries
Vers une Intelligence Artificielle Explicable\\[0.18cm]
pour la Surveillance de l'État Physiologique\\[0.18cm]
de l'Athlète de Haut Niveau
}\\[0.35cm]

{\large\itshape
Espace d'État Fonctionnel, Graphes de Connaissances\\
et Cohérence Sémantique
}

\vspace{0.25cm}
\rule{\linewidth}{0.8pt}

\vspace{0.65cm}

{\large Travail d'Étude et de Recherche}\\[0.10cm]
{\large Année universitaire 2025--2026}

\vspace{0.75cm}

{\large\bfseries TOUHAF Aya}\\[0.45cm]

\renewcommand{\arraystretch}{1.25}
\begin{tabular}{>{\bfseries}ll}
Encadrant & Dr.~Adeel Ahmad \\
Laboratoire & LISIC -- Université du Littoral Côte d'Opale \\
Période & Juin -- Juillet 2026 \\
Date de remise & 14 juillet 2026 \\
\end{tabular}

\vfill
{\small Rapport de Travail d'Étude et de Recherche}

\end{titlepage}

\hypersetup{pageanchor=true}
\pagenumbering{roman}

% Résumé et tables


\section*{Résumé}
\addcontentsline{toc}{section}{Résumé}

\begin{spacing}{1.05}

Les dispositifs portables permettent aujourd'hui d'enregistrer en continu des
signaux physiologiques tels que la fréquence cardiaque (\emph{Heart Rate},
HR)\openaccessref{13}{https://cinc.org/archives/2023/pdf/CinC2023-067.pdf},
la fréquence respiratoire (\emph{Breathing Rate}, BR)\openaccessref{13}{https://cinc.org/archives/2023/pdf/CinC2023-067.pdf}
et les intervalles RR\openaccessref{13}{https://cinc.org/archives/2023/pdf/CinC2023-067.pdf},
c'est-à-dire le temps séparant deux battements cardiaques successifs. Ces
mesures peuvent alimenter des modèles d'apprentissage automatique\openaccessref{6}{https://www.deeplearningbook.org} afin de
reconnaître des états fonctionnels comme le repos, l'effort ou la
récupération. Une classe prédite seule reste toutefois difficile à exploiter :
il faut également comprendre pourquoi le modèle l'a choisie et savoir si cette
décision est suffisamment fiable pour être acceptée sans contrôle humain.

Ce rapport présente notre contribution : une chaîne complète d'intelligence
artificielle explicable (XAI)\openaccessref{1}{https://www.mdpi.com/1099-4300/23/1/18}
appliquée à SportDB~2.0\openaccessref{12}{https://data.mendeley.com/datasets/kzkjkt7mx2/2}. Une première étude est menée
sur la course à pied avec six états fonctionnels. Cinq modèles sont comparés
dans un protocole \emph{Leave-One-Subject-Out} (LOSO)\openaccessref{14}{https://arxiv.org/abs/1904.02666}, c'est-à-dire une
validation dans laquelle un groupe entier est gardé pour le test à chaque
itération : Random Forest (RF, forêt aléatoire)\openaccessref{5}{https://www.stat.berkeley.edu/~breiman/randomforest2001.pdf},
Multilayer Perceptron (MLP, perceptron multicouche)\openaccessref{6}{https://www.deeplearningbook.org},
Long Short-Term Memory (LSTM)\openaccessref{6}{https://www.deeplearningbook.org},
Gated Recurrent Unit (GRU)\openaccessref{7}{https://arxiv.org/abs/1406.1078}
et Transformer\openaccessref{8}{https://arxiv.org/abs/1706.03762}. Les architectures temporelles ne dominent pas
significativement sur l'accuracy globale\openaccessref{6}{https://www.deeplearningbook.org}, c'est-à-dire le taux total de
prédictions correctes, mais elles améliorent le recall\openaccessref{6}{https://www.deeplearningbook.org} de certains états
intermédiaires, autrement dit leur capacité à retrouver les fenêtres qui
appartiennent réellement à ces états.

L'étude est ensuite étendue à quatre disciplines : course, cyclisme, trekking
et patinage. Les décisions du Random Forest sont expliquées avec
\emph{SHapley Additive exPlanations} (SHAP)\openaccessref{3}{https://arxiv.org/abs/1705.07874}, une méthode qui attribue à chaque
variable une contribution positive ou négative à la prédiction. Ces
contributions sont ensuite reliées à des concepts physiologiques au moyen
d'un graphe de connaissances\openaccessref{9}{https://arxiv.org/abs/2003.02320}. Enfin, la confiance prédictive et un score de
cohérence sémantique\openaccessref{21}{https://www.semantic-web-journal.net/system/files/swj2198.pdf} sont combinés dans un mécanisme
\emph{Human-in-the-Loop} (HITL, humain dans la boucle)\openaccessref{22}{https://arxiv.org/abs/2108.00941} qui classe chaque
prédiction en \textsc{accept}, \textsc{check} ou \textsc{review}.

L'évaluation finale est réalisée \emph{out-of-fold}\openaccessref{14}{https://arxiv.org/abs/1904.02666}, ce qui signifie que
chaque fenêtre est prédite et expliquée par un modèle qui n'a pas été entraîné
sur son groupe. Sur 6\,821 fenêtres, 871 prédictions sont classées
\textsc{accept}, 2\,348 \textsc{check} et 3\,602 \textsc{review}. Dans cette
expérimentation, les deux premières catégories présentent les meilleures
accuracies, tandis que \textsc{review} concentre la majorité des erreurs. La
catégorie \textsc{accept} est cependant presque entièrement constituée de
fenêtres TRE, ce qui impose une interprétation prudente. Le système ne cherche
donc pas à remplacer l'expert : il vise à lui montrer l'état prédit, les
variables qui ont soutenu cette prédiction, son niveau de cohérence et les cas
qui nécessitent une vérification.

\noindent\textbf{Mots-clés :} intelligence artificielle explicable, séries
temporelles, surveillance physiologique, SportDB~2.0, Random Forest, LSTM,
GRU, Transformer, SHAP, graphe de connaissances, Leave-One-Subject-Out,
Human-in-the-Loop.

\end{spacing}

\newpage
\tableofcontents

\newpage
\listoffigures
\addcontentsline{toc}{section}{Liste des figures}

\newpage
\listoftables
\addcontentsline{toc}{section}{Liste des tableaux}

\newpage
\pagenumbering{arabic}

% Introduction


\section{Introduction}
\label{sec:introduction}

\subsection{Contexte}

Les montres connectées, les ceintures cardiaques et les capteurs
cardio-respiratoires produisent des séries temporelles, c'est-à-dire des suites
de mesures ordonnées dans le temps. Elles permettent notamment de suivre la
fréquence cardiaque, la respiration et la récupération d'un sportif. Ces
données ouvrent la voie à des systèmes capables de reconnaître
automatiquement un état fonctionnel et d'aider un entraîneur ou un analyste à
mieux comprendre l'évolution d'une séance. Les travaux recensés par Claudino
et al. montrent que l'intelligence artificielle est déjà utilisée dans le
sport pour la prédiction de performance, l'estimation du risque et l'analyse
de la charge d'entraînement~\cite{claudino2019}.

Cette automatisation reste toutefois difficile. Les réponses physiologiques
varient d'un individu à l'autre, mais aussi selon la discipline, l'intensité,
la durée et le capteur utilisé. Une phase d'effort en course ne produit pas
exactement la même dynamique qu'une séance de cyclisme ou de trekking. De
plus, plusieurs fenêtres extraites d'un même enregistrement sont fortement
corrélées. Si certaines de ces fenêtres sont placées dans l'ensemble d'entraînement
(\emph{train}) et d'autres dans l'ensemble de test (\emph{test}), le modèle peut reconnaître des caractéristiques propres au
sujet plutôt que l'état physiologique recherché.

Une deuxième difficulté concerne la compréhension de la décision. Dans un
contexte de surveillance humaine, afficher uniquement le mot « effort » ou
« récupération » n'est pas suffisant. Le lecteur doit pouvoir répondre à
quatre questions simples : quel état a été prédit, avec quelle confiance,
quelles variables ont soutenu cette décision et faut-il demander une
vérification humaine ?

\subsection{Problématique}

La question de recherche retenue est la suivante :

\begin{quote}
\textbf{Comment reconnaître des états physiologiques à partir de données
multi-sport, tout en fournissant une explication sémantique des décisions et
en orientant vers un humain les prédictions insuffisamment fiables ?}
\end{quote}

Cette question se décompose en quatre sous-questions :

\begin{enumerate}[leftmargin=1.8em]
  \item les modèles temporels exploitent-ils mieux les transitions
  physiologiques que les modèles utilisant uniquement des résumés
  statistiques ?
  \item les performances dépendent-elles principalement de l'architecture du
  modèle ou de la séparabilité des données ?
  \item les variables utilisées par le modèle sont-elles cohérentes avec les
  signaux et les profils associés à l'état annoncé ?
  \item une décision combinant confiance prédictive et cohérence sémantique
  permet-elle de sélectionner les prédictions les plus fiables et de montrer
  clairement celles qui doivent être vérifiées ?
\end{enumerate}

\subsection{Objectifs et contributions}

L'objectif n'est pas seulement d'obtenir une classe. Pour chaque fenêtre, le
prototype doit produire une sortie lisible composée de cinq éléments :
l'état prédit, la probabilité associée, les variables les plus influentes, un
score de cohérence sémantique et une recommandation HITL. Cette sortie permet
donc de \emph{voir} la prédiction, mais aussi de comprendre les principaux
arguments utilisés par le modèle.

La contribution principale de ce TER est la conception et l'évaluation d'un
pipeline complet qui relie apprentissage automatique, explicabilité,
représentation sémantique et intervention humaine. Plus précisément, le
travail réalisé comprend :

\begin{itemize}[leftmargin=1.5em]
  \item l'exploration et le prétraitement de SportDB~2.0 ;
  \item la construction d'un espace à six états pour RUN et d'un espace commun
  à trois états pour l'étude multi-sport ;
  \item un fenêtrage commun de 50 échantillons avec un pas de 25, permettant
  de comparer une représentation statistique et une représentation
  séquentielle des mêmes portions de signal ;
  \item la comparaison, dans les mêmes plis LOSO, d'un Random Forest, d'un MLP,
  d'un LSTM, d'un GRU et d'un Transformer ;
  \item l'analyse de la séparabilité des états avec l'analyse de variance,
  le score de Fisher, l'analyse en composantes principales et le score de
  silhouette ;
  \item une étude d'ablation contrôlée sur l'inclusion de SKA ;
  \item l'explication globale et locale des décisions du Random Forest avec
  SHAP ;
  \item la construction d'un graphe de connaissances reliant variables,
  concepts physiologiques et états ;
  \item la définition d'un score de cohérence sémantique ;
  \item la validation out-of-fold d'un mécanisme HITL fondé sur la confiance
  et la cohérence ;
  \item la production de graphiques permettant d'observer les performances,
  les erreurs, les variables importantes et le compromis entre automatisation
  et vérification humaine.
\end{itemize}

\subsection{Organisation du rapport}

La section~\ref{sec:etat_art} présente les travaux existants les plus proches
du sujet et précise la place de notre contribution. La
section~\ref{sec:sujet} décrit notre système, les données retenues et la
construction des états et des variables. La section~\ref{sec:methodologie}
détaille les algorithmes, les formules et le protocole expérimental. La
section~\ref{sec:experiences} présente les résultats et explique la lecture de
chaque graphique. La section~\ref{sec:difficultes} revient sur les difficultés
rencontrées et les solutions mises en place. La
section~\ref{sec:discussion} compare ensuite nos choix et nos résultats aux
travaux étudiés dans l'état de l'art. Enfin, la
section~\ref{sec:conclusion} répond à la problématique et présente les
perspectives.

% État de l'art

\section{État de l'art}
\label{sec:etat_art}

\subsection{Surveillance physiologique par dispositifs portables}

Les dispositifs portables sont utilisés dans le sport pour suivre la charge
interne, la récupération et certains marqueurs de fatigue. Claudino et al.
ont recensé les principales applications de l'intelligence artificielle à la
prédiction de performance et au risque de blessure dans les sports collectifs
~\cite{claudino2019}. Leur synthèse montre que les modèles peuvent exploiter
des données d'entraînement variées, mais souligne aussi la difficulté de
généraliser d'un groupe d'athlètes à un autre.

La base SportDB~2.0, publiée par Romagnoli et al., rassemble des signaux
cardio-respiratoires acquis pendant plusieurs disciplines sportives
~\cite{sportdb2data,sportdb2paper}. Elle contient notamment trois familles de
mesures utilisées dans ce TER :

\begin{itemize}[leftmargin=1.5em]
  \item la fréquence cardiaque, appelée \emph{Heart Rate} (HR), qui indique le
  nombre de battements par minute ;
  \item la fréquence respiratoire, appelée \emph{Breathing Rate} (BR), qui
  indique le nombre de cycles respiratoires par minute ;
  \item les intervalles RR, qui mesurent le temps entre deux pics R
  successifs de l'électrocardiogramme.
\end{itemize}

Lorsque l'intervalle RR est exprimé en millisecondes, la fréquence cardiaque
peut être approximée par :

\begin{equation}
\HR \simeq \frac{60\,000}{\RR}.
\label{eq:hr_rr}
\end{equation}

Dans l'équation~\ref{eq:hr_rr}, $60\,000$ représente le nombre de
millisecondes contenues dans une minute. Si, par exemple, deux battements sont
séparés de 1\,000 millisecondes, alors la fréquence cardiaque estimée est de
$60\,000/1\,000=60$ battements par minute. Cette relation explique pourquoi
HR et RR apportent des informations proches mais inversées : lorsque le
rythme cardiaque augmente, l'intervalle RR diminue généralement.

Des travaux récents se rapprochent de notre problématique. Estrella et
Capdevila utilisent des indicateurs de variabilité de la fréquence cardiaque,
plusieurs modèles d'apprentissage automatique et SHAP afin de distinguer des
profils sportifs~\cite{estrella2025}. Leur objectif porte cependant sur
l'identification de profils ou d'individus, alors que notre travail cherche à
reconnaître des états successifs d'une séance et à décider si chaque
prédiction peut être acceptée ou doit être vérifiée.

\subsection{Reconnaissance d'états et généralisation entre sujets}

Dans un problème de classification supervisée\openaccessref{6}{https://www.deeplearningbook.org}, le modèle apprend une fonction

\begin{equation}
f:\mathcal{X}\rightarrow\mathcal{Y},
\end{equation}

où $\mathcal{X}$ désigne l'espace des observations et $\mathcal{Y}$ l'ensemble
des classes possibles. Dans notre cas, une observation est une fenêtre de
signal et une classe peut être le repos, l'effort ou la récupération.

Les travaux de Dehghani et al. montrent qu'une validation aléatoire peut
surestimer les performances en reconnaissance d'activité humaine lorsque des
fenêtres issues d'un même sujet sont réparties entre le train et le test
~\cite{dehghani2019}. Ils recommandent une validation séparant les sujets.
Cette observation motive le protocole \emph{Leave-One-Subject-Out} (LOSO)
utilisé dans notre travail : à chaque itération, un groupe complet est retiré
de l'entraînement et utilisé pour le test.

Cette exigence est importante pour notre application. Un système destiné à
surveiller un nouvel athlète ne doit pas être évalué uniquement sur des
fenêtres ressemblant à celles qu'il a déjà vues. Le LOSO cherche donc à
mesurer la capacité de généralisation vers un groupe non utilisé pendant
l'apprentissage.

\subsection{Modèles tabulaires et modèles temporels}

Breiman a introduit le Random Forest, ou forêt aléatoire, qui combine de
nombreux arbres de décision afin de réduire la variance d'un arbre unique
~\cite{breiman2001}. Chaque arbre apprend des règles du type « si la fréquence
cardiaque moyenne dépasse un seuil, alors... ». Les prédictions des arbres
sont ensuite agrégées. Ce modèle convient à des variables tabulaires et peut
capturer des relations non linéaires.

Le Multilayer Perceptron (MLP, perceptron multicouche) est un réseau de
neurones composé de couches entièrement connectées. Il reçoit également un
vecteur de variables statistiques, mais apprend des combinaisons plus souples
que les règles d'un arbre.

Les modèles Long Short-Term Memory (LSTM) et Gated Recurrent Unit (GRU) sont
des réseaux récurrents conçus pour traiter des séquences. Le LSTM utilise des
portes qui contrôlent l'information mémorisée, oubliée ou transmise. Le GRU
propose un mécanisme proche avec une structure plus compacte
~\cite{goodfellow2016,cho2014}. Dans notre contexte, ces modèles peuvent
observer l'ordre des valeurs dans une fenêtre, et pas seulement sa moyenne ou
son écart-type.

Le Transformer introduit une autre manière d'analyser les séquences :
l'auto-attention\openaccessref{8}{https://arxiv.org/abs/1706.03762}. Elle est calculée par :

\begin{equation}
\operatorname{Attention}(\mathbf{Q},\mathbf{K},\mathbf{V})
=
\operatorname{softmax}
\left(
\frac{\mathbf{Q}\mathbf{K}^{\top}}{\sqrt{d_k}}
\right)\mathbf{V}.
\label{eq:attention}
\end{equation}

Dans l'équation~\ref{eq:attention}, $\mathbf{Q}$ représente les requêtes,
$\mathbf{K}$ les clés et $\mathbf{V}$ les valeurs. Le produit
$\mathbf{Q}\mathbf{K}^{\top}$ mesure la similarité entre positions de la
séquence. La division par $\sqrt{d_k}$ stabilise les valeurs lorsque la
dimension $d_k$ augmente. La fonction \emph{softmax}\openaccessref{6}{https://www.deeplearningbook.org} transforme ensuite les
scores en poids positifs dont la somme vaut 1. Le modèle peut ainsi accorder
plus d'importance aux instants les plus utiles pour reconnaître l'état.

\subsection{Explicabilité des modèles appliqués aux signaux physiologiques}

Linardatos et al. distinguent les méthodes d'explicabilité globales, qui
décrivent le comportement général d'un modèle, et les méthodes locales, qui
expliquent une prédiction particulière~\cite{linardatos2021}. Ribeiro et al.
ont proposé Local Interpretable Model-agnostic Explanations (LIME)\openaccessref{2}{https://arxiv.org/abs/1602.04938}, qui
approxime localement un modèle complexe par un modèle plus simple.

Lundberg et Lee ont proposé SHapley Additive exPlanations (SHAP), fondé sur
les valeurs de Shapley~\cite{lundberg2017}. Une explication additive s'écrit :

\begin{equation}
g(\mathbf{z}')
=
\phi_0+\sum_{j=1}^{M}\phi_j z'_j.
\label{eq:shap_additive}
\end{equation}

Dans l'équation~\ref{eq:shap_additive}, $\phi_0$ est la valeur de référence du
modèle, $M$ le nombre de variables et $\phi_j$ la contribution de la variable
$j$. Une contribution positive pousse la prédiction vers la classe étudiée ;
une contribution négative l'en éloigne. Le terme $z'_j$ indique si la
variable est présente dans la représentation simplifiée de l'explication.
TreeSHAP\openaccessref{4}{https://arxiv.org/abs/1905.04610} calcule efficacement ces contributions pour les modèles à base
d'arbres.

Shen et al. montrent que l'explicabilité des signaux physiologiques reste un
problème actif : les méthodes doivent expliquer des relations non linéaires,
des dépendances temporelles et des données souvent hétérogènes
~\cite{shen2025}. Dans le domaine sportif, Kranzinger et al. ont recensé 19
études de XAI publiées entre 2014 et juin 2024. Leur analyse montre une forte
domination de SHAP, mais aussi un manque de méthodes sémantiques, de
validation par des experts et de démarches centrées sur l'utilisateur
~\cite{kranzinger2025}. Notre travail se place précisément dans cet espace :
SHAP est utilisé pour attribuer les variables, puis un graphe et un score de
cohérence ajoutent un niveau d'interprétation spécifique au domaine.

\subsection{Graphes de connaissances et intervention humaine}

Un graphe de connaissances représente des informations sous forme d'entités
reliées par des relations. Un triplet peut par exemple être écrit :

\[
(\texttt{mean\_HR},\ \texttt{mesure},\ \texttt{fréquence cardiaque}).
\]

Hogan et al. présentent les principes généraux de ces graphes
~\cite{hogan2020}. Rotmensch et al. montrent qu'ils peuvent aussi structurer
des concepts de santé à partir de données cliniques~\cite{rotmensch2017}.
Lecue souligne leur intérêt pour l'intelligence artificielle explicable :
ils permettent de relier une contribution numérique à un concept et à un
contexte compréhensible par l'utilisateur~\cite{lecue2020}.

Le mécanisme Human-in-the-Loop (HITL, humain dans la boucle) désigne un
système dans lequel l'algorithme ne prend pas toutes les décisions seul. Wu
et al. recensent plusieurs formes de collaboration entre modèle et humain,
notamment l'orientation des cas difficiles vers un expert~\cite{wu2021}.
Geifman et El-Yaniv formalisent cette idée avec la classification sélective\openaccessref{17}{https://arxiv.org/abs/1705.08500} :
le modèle conserve les cas les plus fiables et s'abstient sur les autres.

Deux mesures sont centrales :

\begin{equation}
\text{coverage}
=
\frac{N_{\mathrm{retenu}}}{N_{\mathrm{total}}},
\qquad
\text{risk}
=
1-\text{accuracy}_{\mathrm{retenu}}.
\label{eq:risk_coverage}
\end{equation}

Dans l'équation~\ref{eq:risk_coverage}, la \emph{coverage}\openaccessref{17}{https://arxiv.org/abs/1705.08500} est la proportion
de prédictions que le système accepte de traiter automatiquement. Le
\emph{risk}\openaccessref{17}{https://arxiv.org/abs/1705.08500} est le taux d'erreur parmi ces prédictions. Une bonne stratégie
sélective cherche une coverage suffisamment élevée tout en maintenant un risk
faible.

\subsection{Synthèse des travaux proches et positionnement de notre contribution}

Le tableau~\ref{tab:related_work} résume les travaux les plus proches et met en
évidence la place de notre contribution.

\begin{table}[H]
\centering
\caption{Comparaison synthétique avec des travaux proches.}
\label{tab:related_work}
\begin{tabularx}{\linewidth}{p{2.8cm}XX}
\toprule
\textbf{Travail} & \textbf{Apport principal} & \textbf{Différence avec notre TER} \\
\midrule
Claudino et al.~\cite{claudino2019}
& Revue des usages de l'IA pour la performance et le risque sportif.
& Notre étude construit et évalue un pipeline concret de reconnaissance
d'états physiologiques. \\

Dehghani et al.~\cite{dehghani2019}
& Mise en évidence du biais créé par une validation qui mélange les fenêtres
d'un même sujet.
& Notre protocole reprend cette exigence et l'étend aux prédictions, aux
explications et au mécanisme HITL. \\

Estrella et Capdevila~\cite{estrella2025}
& Classification de profils sportifs avec la variabilité de la fréquence cardiaque
(\emph{Heart Rate Variability}, HRV)\openaccessref{20}{https://www.mdpi.com/2075-4663/13/2/30}, plusieurs modèles et SHAP.
& Notre cible est un état dynamique de séance ; nous ajoutons les modèles
temporels, le graphe sémantique et l'abstention. \\

Kranzinger et al.~\cite{kranzinger2025}
& Revue de l'XAI en science du sport ; domination de SHAP et manque de
validation spécifique au domaine.
& Notre pipeline complète SHAP par un graphe de connaissances et une mesure de
cohérence. \\

Lecue~\cite{lecue2020}
& Rôle des graphes de connaissances pour relier explication et connaissance
symbolique.
& Notre graphe est appliqué aux variables HR, BR et RR et alimente une
décision HITL. \\

Geifman et El-Yaniv~\cite{geifman2017}
& Classification sélective fondée sur la confiance.
& Notre fiabilité combine confiance prédictive et cohérence sémantique. \\
\bottomrule
\end{tabularx}
\end{table}

L'originalité de ce TER ne repose donc pas sur un modèle isolé. Elle vient de
l'intégration, dans une même chaîne, de la comparaison tabulaire/temporelle,
de l'analyse de séparabilité, de SHAP, d'un graphe de connaissances, d'un
score de cohérence et d'une décision HITL évaluée out-of-fold.
\FloatBarrier
% Notre contribution

\section{Notre contribution : conception d'un pipeline explicable de surveillance physiologique multi-sport}
\label{sec:sujet}

\subsection{Vue fonctionnelle et nature des prédictions}

Notre système transforme une portion de signaux cardio-respiratoires en une
décision structurée. Pour une fenêtre donnée, il produit :

\begin{enumerate}[leftmargin=1.8em]
  \item un état fonctionnel prédit, par exemple repos, effort ou récupération ;
  \item une confiance prédictive, comprise entre 0 et 1 ;
  \item les variables qui ont le plus soutenu la classe prédite ;
  \item un score de cohérence entre ces variables et le profil de l'état ;
  \item une recommandation \textsc{accept}, \textsc{check} ou
  \textsc{review}.
\end{enumerate}

La contribution ne se limite donc pas à entraîner un classifieur. Elle ajoute
une représentation de l'explication et un mécanisme de contrôle permettant de
visualiser ce que le modèle a prédit, pourquoi il l'a prédit et si la décision
peut être utilisée automatiquement.

Le prototype vise le suivi de l'état fonctionnel d'un sportif. Il ne réalise
pas de diagnostic médical et ne remplace ni un médecin ni un spécialiste de
la préparation physique.

\subsection{Base de données SportDB~2.0}

SportDB~2.0\cite{sportdb2data} contient 168 jeux de données cardio-respiratoires provenant de
130 sujets et de 11 sports~\cite{sportdb2paper}. Selon les séances, les
fichiers contiennent un électrocardiogramme (ECG)\openaccessref{16}{https://journals.plos.org/plosone/article?id=10.1371/journal.pone.0084018}, la fréquence cardiaque HR,
la fréquence respiratoire BR et les intervalles RR. Des fichiers
\texttt{TrNote.txt} décrivent les phases de la séance et permettent
d'associer une période du signal à un état.

L'arborescence est organisée par discipline et enregistrement. Une séance peut
comprendre :

\begin{itemize}[leftmargin=1.5em]
  \item \texttt{Data.mat}, qui contient les séries numériques ;
  \item \texttt{TrNote.txt}, qui contient les annotations temporelles ;
  \item \texttt{Dem.txt}, qui contient des informations démographiques ou
  contextuelles.
\end{itemize}

\subsection{Choix des disciplines}

La course à pied, notée RUN, a d'abord été étudiée séparément car son protocole
décrit plusieurs phases successives d'effort. Cette granularité permet
d'évaluer si les modèles temporels reconnaissent mieux les transitions.

L'étude multi-sport utilise ensuite :

\begin{itemize}[leftmargin=1.5em]
  \item RUN pour la course ;
  \item CYC pour le cyclisme ;
  \item TRE pour le trekking ;
  \item SKA pour le patinage.
\end{itemize}

Ces disciplines disposent de séries HR et RR directement exploitables dans la
version étudiée. Les sports disponibles principalement sous forme de segments
ECG n'ont pas été intégrés, car leur exploitation aurait nécessité une
détection robuste des pics R, un contrôle des artefacts et une reconstruction
des intervalles RR. Un complexe QRS désigne l'ensemble des ondes Q, R et S
observées sur un électrocardiogramme. Elgendi montre que les méthodes de
détection des complexes QRS\openaccessref{16}{https://journals.plos.org/plosone/article?id=10.1371/journal.pone.0084018} doivent être choisies et validées en fonction du bruit et des contraintes
des systèmes portables~\cite{elgendi2014}.

\subsection{Construction des états fonctionnels}

Pour RUN, six états sont conservés :

\[
\begin{aligned}
S_0&=\text{repos}, &
S_1&=\text{initiation}, &
S_2&=\text{effort intense},\\
S_3&=\text{effort soutenu}, &
S_4&=\text{charge maximale}, &
S_5&=\text{récupération}.
\end{aligned}
\]

Le symbole $S_k$ désigne ici la classe d'indice $k$. Cette représentation
préserve l'ordre des phases du protocole RUN et permet d'étudier les
confusions entre niveaux d'effort voisins.

Pour la comparaison multi-sport, les protocoles sont ramenés à trois états
communs :

\[
S_0=\text{repos},\qquad
S_1=\text{effort},\qquad
S_2=\text{récupération}.
\]

Cette harmonisation est nécessaire parce que les sports ne possèdent pas les
mêmes annotations. Elle permet de construire une cible commune, mais fusionne
les quatre niveaux d'effort détaillés de RUN. La comparaison multi-sport est
donc plus générale, tandis que l'expérience RUN conserve une description plus
fine des transitions.

\subsection{Fenêtrage commun de 50 échantillons avec un pas de 25}

Les signaux bruts sont trop longs pour être fournis en une seule fois aux
modèles. Ils sont découpés en fenêtres de 50 échantillons. Une fenêtre est
simplement une portion continue de 50 valeurs successives du signal.

Le pas est fixé à 25 échantillons. Après avoir extrait une première fenêtre
des positions 1 à 50, la suivante commence à la position 26 et se termine à
75. Les deux fenêtres partagent donc 25 valeurs, ce qui correspond à un
chevauchement de :

\begin{equation}
\text{chevauchement}
=
1-\frac{\text{pas}}{\text{taille de fenêtre}}
=
1-\frac{25}{50}
=
0{,}5
=
50\,\%.
\end{equation}

Ce chevauchement permet de ne pas perdre une transition qui se situerait à la
frontière de deux fenêtres, tout en évitant de créer une nouvelle fenêtre à
chaque échantillon.

Les modèles tabulaires ne reçoivent pas directement les 50 valeurs. Chaque
fenêtre est résumée par neuf variables :

\begin{equation}
\mathbf{x}
=
\left[
\mu_{\HR},
\sigma_{\HR},
\max(\HR),
\mu_{\BR},
\sigma_{\BR},
\max(\BR),
\mu_{\RR},
\sigma_{\RR},
\min(\RR)
\right].
\label{eq:features}
\end{equation}

Le vecteur $\mathbf{x}$ contient donc neuf nombres. Le symbole $\mu$ désigne la
moyenne et $\sigma$ l'écart-type, qui mesure la dispersion des valeurs autour
de la moyenne. Plus précisément :

\begin{itemize}[leftmargin=1.5em]
  \item $\mu_{\HR}$ est la fréquence cardiaque moyenne de la fenêtre ;
  \item $\sigma_{\HR}$ mesure sa variabilité ;
  \item $\max(\HR)$ est la plus grande valeur de HR ;
  \item les trois mêmes résumés sont calculés pour BR ;
  \item $\mu_{\RR}$ et $\sigma_{\RR}$ décrivent les intervalles RR ;
  \item $\min(\RR)$ correspond à l'intervalle le plus court, donc souvent au
  rythme cardiaque le plus rapide de la fenêtre.
\end{itemize}

Le Random Forest et le MLP utilisent ce vecteur de neuf variables. Les modèles
LSTM, GRU et Transformer conservent au contraire la séquence complète de
50 échantillons afin d'observer l'ordre et l'évolution des signaux. Les deux
familles de modèles analysent ainsi les mêmes portions temporelles, mais sous
deux représentations différentes.

Les fenêtres contenant des valeurs non finies sont exclues. La normalisation
est apprise uniquement sur les données d'entraînement de chaque pli afin
d'éviter que le groupe de test influence le modèle.

\subsection{Corpus expérimentaux retenus}

Avec le fenêtrage de 50 échantillons et le pas de 25, le corpus RUN utilisé
pour la comparaison des cinq modèles contient 1\,251 fenêtres réparties entre
dix groupes et six états. Le corpus multi-sport contient 6\,821 fenêtres et
32 groupes dans la version expérimentale évaluée.

\begin{table}[H]
\centering
\caption{Composition du corpus multi-sport utilisé dans les expériences.}
\label{tab:corpus_multisport}
\begin{tabular}{lrrr}
\toprule
\textbf{Sport} & \textbf{Groupes} & \textbf{Fenêtres} & \textbf{Part} \\
\midrule
RUN & 10 & 1\,296 & 19,0\,\% \\
CYC & 11 & 1\,091 & 16,0\,\% \\
TRE & 7  & 2\,278 & 33,4\,\% \\
SKA & 4  & 2\,156 & 31,6\,\% \\
\midrule
\textbf{Total} & \textbf{32} & \textbf{6\,821} & \textbf{100\,\%} \\
\bottomrule
\end{tabular}
\end{table}

La dernière colonne se lit comme la proportion du corpus total. Par exemple,
les 2\,278 fenêtres TRE représentent 33,4\,\% des 6\,821 fenêtres. La
répartition des classes est déséquilibrée : 1\,275 fenêtres de repos,
4\,796 fenêtres d'effort et 750 fenêtres de récupération. TRE ne contribue
qu'à l'état effort dans l'espace commun. Cette particularité doit être prise
en compte lors de l'interprétation des accuracies et du mécanisme HITL.

\subsection{Organisation temporelle du travail}
\label{subsec:organisation_temporelle}

Le projet a officiellement commencé le 1\textsuperscript{er} juin 2026,
à la suite d'une première réunion avec l'encadrant. Le travail a ensuite
progressé de manière itérative : les recherches bibliographiques, les
expérimentations et les choix méthodologiques ont été approfondis au fur et à
mesure des résultats obtenus et des réunions de suivi.

La figure~\ref{fig:gantt_ter} présente cette progression sous la forme d'un
\textbf{diagramme de Gantt}. Les losanges correspondent aux réunions, aux présentations ou aux
jalons importants. Les dates indiquées reposent sur les réunions
d'encadrement, l'historique Git et les dates de modification des principaux
scripts.

\begin{landscape}
\begin{figure}[p]
\centering
\scriptsize

\begin{ganttchart}[
    x unit=0.34cm,
    y unit title=0.62cm,
    y unit chart=0.49cm,
    title height=1,
    bar height=0.38,
    milestone height=0.38,
    group height=0.26,
    hgrid,
    vgrid,
    title label font=\bfseries\tiny,
    bar label font=\tiny,
    milestone label font=\tiny,
    group label font=\bfseries\tiny,
    bar label node/.append style={align=right,text width=6.3cm},
    milestone label node/.append style={align=right,text width=6.3cm},
    group label node/.append style={align=right,text width=6.3cm},
    bar/.append style={fill=blue!45},
    group/.append style={fill=blue!70},
    milestone/.append style={fill=orange!80}
]{1}{44}

\gantttitle{Juin 2026}{30}
\gantttitle{Juillet 2026}{14} \\

\gantttitle{1--7 juin}{7}
\gantttitle{8--14 juin}{7}
\gantttitle{15--21 juin}{7}
\gantttitle{22--28 juin}{7}
\gantttitle{29 juin--5 juillet}{7}
\gantttitle{6--12 juillet}{7}
\gantttitle{13--14 juillet}{2} \\

\ganttgroup
{Phase 1 : compréhension du sujet et premier prototype}
{1}{18} \\

\ganttmilestone
{Réunion de lancement et définition du sujet}
{1} \\

\ganttbar
{Recherche d'informations, lecture d'articles et compréhension du sujet}
{1}{11} \\

\ganttbar
{Exploration de SportDB~2.0 et premiers développements}
{4}{11} \\

\ganttbar
{Premier prototype Random Forest appliqué uniquement à RUN}
{7}{12} \\

\ganttmilestone
{Présentation du premier résultat RUN et validation de l'orientation}
{12} \\

\ganttgroup
{Phase 2 : comparaison des modèles sur RUN}
{18}{28} \\

\ganttmilestone
{Réunion de suivi : comparaison des modèles puis extension multi-sport}
{18} \\

\ganttmilestone
{Création du dépôt Git et intégration du prototype}
{19} \\

\ganttbar
{Structuration du pipeline de modélisation des états fonctionnels}
{19}{20} \\

\ganttbar
{Comparaison Random Forest, MLP, LSTM, GRU et Transformer sur RUN}
{20}{27} \\

\ganttbar
{Analyse des résultats RUN et du recall des différents états}
{24}{28} \\

\ganttbar
{Recherche bibliographique complémentaire selon les besoins}
{19}{28} \\

\ganttgroup
{Phase 3 : extension multi-sport et analyse des données}
{29}{34} \\

\ganttbar
{Construction du corpus commun RUN, CYC, TRE et SKA}
{29}{30} \\

\ganttbar
{Comparaison des cinq modèles dans l'expérience multi-sport}
{29}{32} \\

\ganttbar
{Étude d'ablation sur l'inclusion de SKA}
{29}{31} \\

\ganttbar
{Analyse des distributions et projection en composantes principales}
{29}{32} \\

\ganttbar
{Diagnostic multi-sport et analyse de la séparabilité}
{32}{34} \\

\ganttgroup
{Phase 4 : explicabilité, graphe de connaissances et HITL}
{40}{41} \\

\ganttbar
{Validation des explications et comparaison entre SHAP et ANOVA}
{40}{40} \\

\ganttbar
{Construction du graphe et calcul de la cohérence sémantique}
{40}{40} \\

\ganttbar
{Mise en place du mécanisme Human-in-the-Loop}
{40}{41} \\

\ganttbar
{Évaluation out-of-fold des explications et de la fiabilité}
{41}{41} \\

\ganttgroup
{Phase 5 : rédaction et finalisation du rapport}
{41}{44} \\

\ganttbar
{Rédaction du rapport et organisation des contributions}
{41}{44} \\

\ganttbar
{Explication des résultats, des formules et des figures}
{41}{44} \\

\ganttbar
{Vérification des scripts, des références et du document}
{42}{44} \\

\ganttmilestone
{Finalisation et date de remise du rapport}
{44}

\end{ganttchart}

\caption{Diagramme de Gantt présentant l'organisation chronologique du TER
entre le 1\textsuperscript{er} juin et le 14 juillet 2026.}
\label{fig:gantt_ter}
\end{figure}
\end{landscape}

Le déroulement du projet montre une progression en cinq phases. La première
phase a été consacrée à la compréhension du sujet, à la consultation de la
littérature scientifique et à la création d'un premier prototype sur RUN.
La présentation du 12 juin a permis de vérifier la faisabilité de la
classification des états physiologiques avec un Random Forest.

La réunion du 18 juin a ensuite conduit à élargir l'étude. Le travail a alors
porté sur la comparaison de cinq familles de modèles : Random Forest,
Multilayer Perceptron, Long Short-Term Memory, Gated Recurrent Unit et
Transformer. Cette comparaison a d'abord été réalisée sur RUN, avant d'être
étendue à plusieurs sports.

À partir du 29 juin, la contribution a été complétée par la construction du
corpus multi-sport, l'étude d'ablation, l'analyse des distributions et les
mesures de séparabilité. Ces expériences ont permis d'étudier l'influence de
l'architecture du modèle ainsi que celle de la structure des données.

Le 10 juillet, le travail a été poursuivi avec la partie explicabilité. Les
prédictions du Random Forest ont été étudiées avec SHAP, puis reliées à un
graphe de connaissances. La confiance prédictive et la cohérence sémantique
ont enfin été combinées dans un mécanisme Human-in-the-Loop.

La dernière phase, du 11 au 14 juillet, a été consacrée à la rédaction, à
l'explication détaillée de la méthode et des résultats, ainsi qu'à la
vérification de la cohérence générale du rapport. Cette chronologie montre que
la contribution finale résulte d'une évolution progressive, depuis un premier
classifieur appliqué à RUN jusqu'à une chaîne complète permettant de prédire,
d'expliquer et d'orienter les décisions vers un contrôle humain.


% Méthodologie

\section{Méthodologie et plan expérimental}
\label{sec:methodologie}

\subsection{Chaîne de traitement complète}

La figure~\ref{fig:pipeline} résume notre pipeline. Elle se lit de gauche à
droite. Le premier bloc contient les signaux bruts HR, BR et RR. Le deuxième
bloc réalise le fenêtrage et construit soit un vecteur statistique, soit une
séquence. Le troisième bloc correspond aux cinq modèles de classification. Le
quatrième bloc explique la décision du Random Forest avec SHAP. Le cinquième
bloc relie les variables importantes aux concepts physiologiques du graphe de
connaissances. Le dernier bloc combine confiance et cohérence pour produire la
décision HITL.

\begin{figure}[H]
\centering
\safeincludegraphics[width=0.96\linewidth]
{outputs_validation/4_pipeline_summary.png}
\caption{Chaîne complète proposée : acquisition des signaux, fenêtrage,
classification, explication SHAP, mise en relation sémantique et décision HITL.}
\label{fig:pipeline}
\end{figure}

Cette figure met en évidence la contribution centrale du TER : les résultats
du classifieur ne sont pas considérés comme une fin. Ils sont transmis à deux
modules supplémentaires, l'explication et la qualification de fiabilité, afin
de produire une sortie compréhensible et contrôlable.

\subsection{Prétraitement et normalisation}

Le chargement parcourt les dossiers de sport, de groupe et de séance. Les
tableaux MATLAB sont convertis en nombres flottants. Les annotations du fichier
\texttt{TrNote.txt} permettent ensuite d'associer chaque portion du signal à
un état fonctionnel.

Les fichiers incomplets, les fenêtres trop courtes et les valeurs non finies
sont exclus. Pour empêcher une variable numériquement grande de dominer les
autres, les données sont standardisées\openaccessref{6}{https://www.deeplearningbook.org} :

\begin{equation}
x'_{i,j}
=
\frac{x_{i,j}-\mu^{\mathrm{train}}_j}
{\sigma^{\mathrm{train}}_j}.
\label{eq:standardization}
\end{equation}

Dans l'équation~\ref{eq:standardization}, $x_{i,j}$ est la valeur de la
variable $j$ pour la fenêtre $i$. La moyenne
$\mu^{\mathrm{train}}_j$ et l'écart-type
$\sigma^{\mathrm{train}}_j$ sont calculés uniquement sur le train. La valeur
standardisée $x'_{i,j}$ indique alors combien d'écarts-types la mesure se
situe au-dessus ou au-dessous de la moyenne d'entraînement. Le groupe de test
n'intervient jamais dans le calcul de l'outil de standardisation, appelé
\emph{scaler} en apprentissage automatique.

\subsection{Modèles comparés}

Les cinq modèles sont évalués sur les mêmes plis. Le
tableau~\ref{tab:model_config} résume leur rôle.

\begin{table}[H]
\centering
\caption{Modèles utilisés et information exploitée par chacun.}
\label{tab:model_config}
\begin{tabularx}{\linewidth}{lXX}
\toprule
\textbf{Modèle} & \textbf{Entrée} & \textbf{Principe} \\
\midrule
Random Forest
& 9 variables statistiques
& Agrégation de 200 à 250 arbres selon l'expérience, profondeur maximale 8,
minimum 3 observations par feuille et pondération des classes. \\

Multilayer Perceptron
& 9 variables statistiques
& Réseau dense qui apprend des combinaisons non linéaires des neuf variables,
avec activation \emph{Rectified Linear Unit} (ReLU, unité linéaire rectifiée)\openaccessref{6}{https://www.deeplearningbook.org}, sortie softmax et arrêt anticipé. \\

Long Short-Term Memory
& Séquences de 50 échantillons
& Réseau récurrent à mémoire qui suit l'évolution du signal dans la fenêtre. \\

Gated Recurrent Unit
& Séquences de 50 échantillons
& Réseau récurrent compact utilisant des portes de mise à jour et de
réinitialisation. \\

Transformer
& Séquences de 50 échantillons
& Modèle d'auto-attention qui compare directement les positions de la
séquence et utilise un encodage de leur ordre. \\
\bottomrule
\end{tabularx}
\end{table}

La sortie \emph{softmax} transforme les scores du réseau en probabilités dont
la somme vaut 1. La classe ayant la probabilité la plus élevée devient la
classe prédite. L'arrêt anticipé, ou \emph{early stopping}\openaccessref{6}{https://www.deeplearningbook.org}, interrompt
l'entraînement lorsqu'aucune amélioration suffisante n'est observée, afin de
limiter le surapprentissage et le temps de calcul.

\subsection{Validation Leave-One-Subject-Out}

Pour un groupe $g$, l'ensemble de test est défini par :

\begin{equation}
\mathcal{D}^{(g)}_{\mathrm{test}}
=
\left\{
(\mathbf{x}_i,y_i)
\mid
group_i=g
\right\}.
\label{eq:loso}
\end{equation}

Dans l'équation~\ref{eq:loso}, $\mathbf{x}_i$ est l'entrée de la fenêtre $i$,
$y_i$ sa vraie classe et $group_i$ son groupe. Toutes les fenêtres du groupe
$g$ sont placées dans le test. Les autres groupes forment le train. La
procédure est répétée jusqu'à ce que chaque groupe ait été utilisé une fois
comme test.

Cette validation réduit le risque de fuite entre fenêtres voisines et permet
d'estimer la capacité du modèle à généraliser à un groupe non observé pendant
l'entraînement.

\subsection{Métriques de classification}

L'accuracy mesure la proportion totale de prédictions correctes :

\begin{equation}
\text{accuracy}
=
\frac{N_{\mathrm{correct}}}{N_{\mathrm{total}}},
\label{eq:accuracy}
\end{equation}

où $N_{\mathrm{correct}}$ est le nombre de fenêtres bien classées et
$N_{\mathrm{total}}$ le nombre de fenêtres testées.

Le recall d'une classe $k$ mesure la part des vraies fenêtres de cette classe
qui ont été retrouvées :

\begin{equation}
\text{recall}_k
=
\frac{TP_k}{TP_k+FN_k}.
\label{eq:recall}
\end{equation}

$TP_k$, pour \emph{True Positives}, désigne les fenêtres de la classe $k$
correctement reconnues. $FN_k$, pour \emph{False Negatives}, désigne les
fenêtres de cette classe attribuées par erreur à une autre classe. Un recall
élevé est donc important lorsqu'on cherche à ne pas manquer un état.

La matrice de confusion place les vraies classes sur un axe et les classes
prédites sur l'autre. Les valeurs de la diagonale correspondent aux bonnes
prédictions ; les cases hors diagonale indiquent les confusions.

Deux agrégations de l'accuracy sont distinguées :

\begin{itemize}[leftmargin=1.5em]
  \item la moyenne des accuracies des plis, où chaque groupe compte autant ;
  \item l'accuracy poolée, où chaque fenêtre compte autant.
\end{itemize}

Ces valeurs peuvent être différentes lorsqu'un groupe contient beaucoup plus
de fenêtres qu'un autre.

\subsection{Tests statistiques}

Le test de Friedman\openaccessref{15}{https://www.jmlr.org/papers/volume7/demsar06a/demsar06a.pdf} compare les rangs des cinq modèles sur les mêmes groupes.
Son hypothèse nulle affirme qu'aucun modèle ne se distingue globalement. Une
$p$-value supérieure à 0,05 signifie que les différences observées ne
permettent pas de rejeter cette hypothèse au seuil de 5\,\%.

Le test de Wilcoxon\openaccessref{15}{https://www.jmlr.org/papers/volume7/demsar06a/demsar06a.pdf} compare deux conditions appariées, ici l'entraînement avec
et sans SKA. Il tient compte du fait que les deux résultats proviennent des
mêmes groupes de test. Demšar recommande ces tests non paramétriques pour
comparer plusieurs classifieurs sur des ensembles appariés
~\cite{demsar2006}.

\subsection{Analyse de séparabilité}

L'analyse de variance, ou \emph{Analysis of Variance} (ANOVA)\openaccessref{26}{https://www.itl.nist.gov/div898/handbook/prc/section4/prc431.htm}, vérifie si la
moyenne d'une variable diffère entre les classes. Une statistique $F$ élevée
indique que la variance entre classes est grande par rapport à la variance à
l'intérieur des classes.

Le score de Fisher\openaccessref{25}{https://arxiv.org/abs/1202.3725} utilisé pour une variable peut être résumé par :

\begin{equation}
F_{\mathrm{Fisher}}
=
\frac{\text{variance entre les classes}}
{\text{variance à l'intérieur des classes}}.
\label{eq:fisher}
\end{equation}

Une valeur élevée signifie que les centres des classes sont éloignés par
rapport à leur dispersion.

L'analyse en composantes principales, ou \emph{Principal Component Analysis}
(PCA)\openaccessref{23}{https://arxiv.org/abs/1404.1100}, transforme les neuf variables en nouveaux axes. Le premier axe explique
la plus grande part possible de la variance, puis le deuxième explique la
plus grande part restante tout en étant perpendiculaire au premier. La PCA
permet donc de projeter les fenêtres en deux dimensions pour observer les
regroupements.

Le score de silhouette\openaccessref{24}{https://wis.kuleuven.be/stat/robust/papers/publications-1987/rousseeuw-silhouettes-jcam-sciencedirectopenarchiv.pdf} d'une fenêtre $i$ est :

\begin{equation}
s(i)
=
\frac{b(i)-a(i)}
{\max\{a(i),b(i)\}},
\label{eq:silhouette}
\end{equation}

où $a(i)$ est la distance moyenne entre la fenêtre et les fenêtres de sa
propre classe, et $b(i)$ la plus petite distance moyenne vers une autre
classe. Un score proche de 1 indique une bonne séparation ; un score proche de
0 indique un chevauchement ; un score négatif signifie que la fenêtre est en
moyenne plus proche d'une autre classe.

\subsection{Explications SHAP et graphe de connaissances}

Le Random Forest est retenu pour la partie explicabilité parce qu'il obtient
la meilleure moyenne LOSO multi-sport dans l'expérimentation et qu'il est
compatible avec TreeSHAP. Pour une fenêtre, SHAP fournit une contribution
$\phi_j$ pour chaque variable. Les contributions positives vers la classe
prédite sont considérées comme des arguments de soutien.

Le graphe de connaissances relie trois niveaux :

\begin{enumerate}[leftmargin=1.8em]
  \item les variables calculées, comme \texttt{mean\_HR} ;
  \item les concepts physiologiques HR, BR et RR ;
  \item les états fonctionnels repos, effort et récupération.
\end{enumerate}

\begin{figure}[H]
\centering
\safeincludegraphics[width=0.96\linewidth]
{outputs_explainability_kg/3_knowledge_graph.png}
\caption{Graphe de connaissances utilisé pour traduire les variables
statistiques en concepts physiologiques puis en états fonctionnels.}
\label{fig:knowledge_graph}
\end{figure}

La figure~\ref{fig:knowledge_graph} se lit en suivant les arêtes entre les
nœuds. Les nœuds représentant les variables sont reliés au signal dont elles
proviennent. Les concepts HR, BR et RR sont ensuite reliés aux états auxquels
ils peuvent contribuer. Le graphe n'apprend pas la prédiction à la place du
modèle ; il organise l'explication dans un vocabulaire plus compréhensible.

\subsection{Cohérence sémantique}

Pour chaque sport et chaque état, des profils de référence sont calculés à
partir des médianes du train. La médiane est la valeur qui partage les
observations en deux moitiés et résiste mieux aux valeurs extrêmes qu'une
moyenne.

La cohérence sémantique est définie par :

\begin{equation}
C(\mathbf{x},\hat{y})
=
\frac{
\sum_{j\in P(\mathbf{x},\hat{y})}
|\phi_j|\,c_j(\mathbf{x},\hat{y})
}{
\sum_{j\in P(\mathbf{x},\hat{y})}
|\phi_j|
}.
\label{eq:consistency}
\end{equation}

Dans cette formule :

\begin{itemize}[leftmargin=1.5em]
  \item $\mathbf{x}$ est la fenêtre analysée ;
  \item $\hat{y}$ est la classe prédite ;
  \item $P(\mathbf{x},\hat{y})$ contient au plus les cinq variables qui
  soutiennent le plus cette classe ;
  \item $|\phi_j|$ est l'importance absolue de la contribution SHAP ;
  \item $c_j(\mathbf{x},\hat{y})$ est la similarité entre la valeur observée
  et le profil de l'état prédit.
\end{itemize}

La formule réalise une moyenne pondérée. Une variable très importante selon
SHAP compte davantage qu'une variable secondaire. Le score final est compris
entre 0 et 1 : une valeur élevée signifie que les arguments principaux du
modèle sont proches du profil appris pour l'état annoncé.

Lorsqu'un sport ne contient pas les trois états, des profils globaux calculés
sur le train sont utilisés. Cette précaution évite qu'un sport ne possédant
qu'une seule classe obtienne automatiquement une cohérence parfaite.

\subsection{Fiabilité et décisions Human-in-the-Loop}

La confiance du modèle pour la classe choisie est notée $p_i$ et la cohérence
sémantique $C_i$. La fiabilité est définie par :

\begin{equation}
R_i=\sqrt{p_iC_i}.
\label{eq:reliability}
\end{equation}

Le produit $p_iC_i$ impose que les deux composantes soient élevées. Une forte
confiance ne suffit pas si l'explication est incohérente, et une bonne
cohérence ne suffit pas si le modèle hésite. La racine carrée maintient le
score dans l'intervalle $[0,1]$ et correspond à une moyenne géométrique.

Les catégories sont définies par :

\begin{itemize}[leftmargin=1.5em]
  \item \textsc{accept} si $p_i\geq 0{,}75$ et $C_i\geq 0{,}70$ : la
  prédiction est suffisamment confiante et cohérente pour être placée dans le
  groupe le plus fiable ;
  \item \textsc{check} si $p_i\geq 0{,}55$ et $C_i\geq 0{,}45$ : la
  prédiction paraît exploitable mais mérite un contrôle rapide ;
  \item \textsc{review} dans les autres cas : la décision est orientée vers
  une vérification humaine complète.
\end{itemize}

Ces seuils sont expérimentaux. Ils ne constituent ni une norme clinique ni
une garantie générale.

\subsection{Explicabilité entièrement out-of-fold}

Dans l'évaluation finale, le modèle, le scaler, les profils, les valeurs SHAP
et la décision HITL sont recalculés dans chaque pli LOSO. Une fenêtre du groupe
test est donc prédite et expliquée sans avoir servi à construire le modèle ou
les profils de référence.

Cette étape est une contribution méthodologique importante du TER. Elle évite
de mesurer la qualité des explications sur des données déjà utilisées pendant
l'entraînement et rend la comparaison entre confiance, cohérence et erreur
plus réaliste.

% Expériences

\section{Expériences, résultats et interprétation}
\label{sec:experiences}

\subsection{Reconnaissance détaillée des six états de RUN}

Cette première expérience étudie la course à pied avec six états. Les cinq
modèles sont entraînés et testés sur les mêmes dix plis LOSO. L'objectif est
de vérifier si l'information temporelle aide à distinguer des phases d'effort
proches.

\begin{figure}[H]
\centering
\safeincludegraphics[width=0.78\linewidth]
{outputs_dl/confusion_RandomForest.png}
\caption{Matrice de confusion LOSO du Random Forest sur les six états de RUN.}
\label{fig:run_confusion}
\end{figure}

Pour lire la figure~\ref{fig:run_confusion}, il faut comparer la vraie classe
et la classe prédite. Les cases de la diagonale correspondent aux fenêtres
correctement classées. Une case située hors de la diagonale indique qu'un état
a été confondu avec un autre. Plus la couleur est intense, plus le nombre ou
la proportion de fenêtres concernées est important.

La diagonale est plus marquée pour les états les plus distincts, tandis que
les phases d'effort voisines se mélangent davantage. Cette observation est
cohérente avec le protocole : l'initiation, l'effort intense, l'effort soutenu
et la charge maximale appartiennent à une même séance et peuvent partager des
niveaux de HR, BR et RR proches. La matrice permet donc de voir précisément
où l'accuracy globale est perdue.

\begin{table}[H]
\centering
\caption{Accuracy LOSO sur RUN à six états.}
\label{tab:run_models}
\begin{tabular}{lcc}
\toprule
\textbf{Modèle} & \textbf{Accuracy moyenne} & \textbf{Écart-type} \\
\midrule
Random Forest & 0,554 & 0,097 \\
MLP           & 0,555 & 0,071 \\
LSTM          & 0,547 & 0,085 \\
GRU           & \textbf{0,584} & 0,125 \\
Transformer   & 0,549 & 0,114 \\
\bottomrule
\end{tabular}
\end{table}

L'accuracy moyenne du GRU est la plus élevée, avec 0,584, soit environ
58,4\,\% de prédictions correctes en moyenne par pli. L'écart-type de 0,125
montre cependant que les performances changent fortement d'un groupe à
l'autre. Le test de Friedman donne $\chi^2=2{,}40$ et $p=0{,}66$. Comme la
$p$-value est largement supérieure à 0,05, les écarts ne permettent pas de
conclure qu'une architecture domine globalement les autres.

L'accuracy seule peut masquer un avantage limité à une classe. Le
tableau~\ref{tab:run_s2_recall} présente donc le recall de l'état $S_2$.

\begin{table}[H]
\centering
\caption{Recall de l'état $S_2$, correspondant à l'effort intense.}
\label{tab:run_s2_recall}
\begin{tabular}{lc}
\toprule
\textbf{Modèle} & \textbf{Recall de $S_2$} \\
\midrule
Random Forest & 30,2\,\% \\
MLP           & 16,5\,\% \\
LSTM          & 52,8\,\% \\
GRU           & 52,8\,\% \\
Transformer   & \textbf{54,2\,\%} \\
\bottomrule
\end{tabular}
\end{table}

Le Transformer retrouve 54,2\,\% des vraies fenêtres $S_2$, contre 30,2\,\%
pour le Random Forest et 16,5\,\% pour le MLP. Les trois modèles séquentiels
dépassent ici les modèles tabulaires. Cela suggère que l'ordre des valeurs
dans les 50 échantillons contient une information utile pour reconnaître
l'effort intense, même si ce gain local ne suffit pas à améliorer
significativement l'accuracy totale.

\begin{figure}[H]
\centering
\safeincludegraphics[width=0.90\linewidth]
{outputs_dl/1_model_comparison.png}
\caption{Comparaison des cinq modèles sur RUN ; les barres résument leur
performance moyenne et leur variabilité entre groupes.}
\label{fig:run_model_comparison}
\end{figure}

Dans la figure~\ref{fig:run_model_comparison}, la hauteur de chaque barre
représente la performance moyenne. Les barres d'erreur représentent la
variabilité entre plis. Les intervalles se chevauchent fortement, ce qui
explique pourquoi le test de Friedman n'identifie pas de différence globale
significative. Le graphique doit donc être lu avec le tableau des recalls :
la moyenne générale est proche, mais le comportement par classe diffère.

\subsection{Comparaison multi-sport}

L'expérience multi-sport utilise les états repos, effort et récupération et
32 groupes LOSO. Le tableau~\ref{tab:multisport_models} présente, pour chaque
modèle, la moyenne, l'écart-type, le minimum et le maximum observés entre les
groupes.

\begin{table}[H]
\centering
\caption{Résultats LOSO moyens sur le corpus multi-sport.}
\label{tab:multisport_models}
\begin{tabular}{lrrrr}
\toprule
\textbf{Modèle} &
\textbf{Moyenne} &
\textbf{Écart-type} &
\textbf{Minimum} &
\textbf{Maximum} \\
\midrule
Random Forest & \textbf{0,7995} & 0,1838 & 0,3973 & 1,0000 \\
MLP           & 0,7827 & 0,1905 & 0,3868 & 1,0000 \\
LSTM          & 0,7105 & 0,2268 & 0,2049 & 1,0000 \\
GRU           & 0,7509 & 0,2068 & 0,2643 & 1,0000 \\
Transformer   & 0,7645 & 0,1909 & 0,3252 & 0,9873 \\
\bottomrule
\end{tabular}
\end{table}

Le Random Forest obtient la meilleure moyenne, 0,7995. Cette valeur ne
signifie pas que 79,95\,\% de toutes les fenêtres sont correctes : elle est la
moyenne non pondérée des 32 accuracies de plis. Chaque groupe compte donc
autant, quelle que soit sa taille.

Les minimums proches de 0,20 à 0,40 et les maximums proches de 1 montrent une
forte hétérogénéité. Certains groupes sont presque entièrement reconnus,
d'autres sont très difficiles. Le test de Friedman donne
$\chi^2=4{,}15$ et $p=0{,}39$ ; aucune supériorité globale n'est donc
démontrée au seuil de 5\,\%.

\begin{figure}[H]
\centering
\safeincludegraphics[width=0.88\linewidth]
{outputs_dl_multisport/1_model_comparison_multisport.png}
\caption{Accuracy LOSO moyenne des cinq architectures sur le corpus
multi-sport.}
\label{fig:multisport_comparison}
\end{figure}

La figure~\ref{fig:multisport_comparison} donne une vue immédiate du classement
moyen. Le Random Forest et le MLP apparaissent légèrement au-dessus, mais les
barres d'erreur sont grandes et se recouvrent. Le graphique confirme donc que
la variabilité entre groupes est plus forte que la différence moyenne entre
architectures.

\begin{figure}[H]
\centering
\safeincludegraphics[width=0.94\linewidth]
{outputs_dl_multisport/4_per_sport_accuracy.png}
\caption{Accuracy des modèles selon la discipline du groupe testé.}
\label{fig:per_sport_accuracy}
\end{figure}

Dans la figure~\ref{fig:per_sport_accuracy}, chaque groupe de barres correspond
à un sport. La comparaison verticale indique le niveau de performance dans ce
sport ; la comparaison horizontale montre comment les modèles changent d'une
discipline à l'autre. RUN et CYC sont souvent mieux reconnus. Dans les données
préparées pour cette expérimentation, SKA reste proche de 0,40 pour plusieurs
modèles, tandis que TRE présente des groupes très faciles et d'autres plus
difficiles.

Le Transformer obtient le meilleur recall de la récupération, avec
35,7\,\%, devant le GRU à 33,2\,\%, le Random Forest à 29,7\,\%, le LSTM à
28,8\,\% et le MLP à 12,1\,\%. Comme dans RUN, l'information temporelle semble
donc utile pour une classe précise sans produire de domination globale.

Deux accuracies doivent être distinguées :

\begin{itemize}[leftmargin=1.5em]
  \item 0,7995 est la moyenne des 32 accuracies de plis ;
  \item 0,6995 est l'accuracy poolée sur les 6\,821 prédictions out-of-fold.
\end{itemize}

L'accuracy poolée se calcule en réunissant toutes les fenêtres de test. Les
grands groupes pèsent alors davantage. Les groupes SKA contiennent beaucoup de
fenêtres et obtiennent de faibles performances dans cette version ; ils
abaissent donc davantage la valeur poolée que la moyenne par pli.

\subsection{Analyse de séparabilité}

L'analyse de séparabilité cherche à savoir si les classes sont déjà distinctes
dans les variables d'entrée, avant même de choisir un modèle complexe.

\begin{figure}[H]
\centering
\safeincludegraphics[width=0.96\linewidth]
{outputs_distribution_analysis/2_pca_projection.png}
\caption{Projection PCA des fenêtres ; chaque point représente une fenêtre et
les deux axes résument 73,4\,\% de la variance totale.}
\label{fig:pca}
\end{figure}

Dans la figure~\ref{fig:pca}, chaque point est une fenêtre décrite par neuf
variables. Les axes horizontaux et verticaux ne sont pas des signaux
physiologiques directs : ce sont deux combinaisons linéaires produites par la
PCA. Des points proches possèdent des profils statistiques proches. La forme
ou la couleur permet de distinguer l'état et le sport.

Les deux premières composantes expliquent 73,4\,\% de la variance. La
projection conserve donc une grande partie, mais pas la totalité, de
l'information. Les regroupements apparaissent fortement structurés par sport,
ce qui montre que la discipline influence la distribution des variables. La
PCA est une visualisation : elle suggère des tendances, mais ne prouve pas à
elle seule qu'une classe est prédictible.

\begin{table}[H]
\centering
\caption{Indicateurs de séparabilité calculés par sport.}
\label{tab:separability}
\begin{tabular}{lrrrl}
\toprule
\textbf{Sport} &
\textbf{Silhouette} &
\textbf{Fisher HR} &
\textbf{ANOVA mean\_HR} &
\textbf{Interprétation} \\
\midrule
RUN & $+0,439$ & 7,445 & 3264,33 & séparation nette \\
CYC & $+0,174$ & 2,313 & 243,88  & séparation modérée \\
SKA & $-0,051$ & 0,014 & 17,76   & fort chevauchement \\
TRE & N/A      & N/A   & N/A     & un seul état \\
\bottomrule
\end{tabular}
\end{table}

RUN possède le meilleur score de silhouette et le plus grand score de Fisher
sur HR. Cela signifie que ses états sont relativement éloignés par rapport à
leur dispersion interne. CYC présente une séparation plus faible mais encore
positive. Dans la version expérimentale étudiée, SKA obtient une silhouette
légèrement négative et un Fisher HR proche de zéro, ce qui indique un
chevauchement important des classes dans les variables disponibles. Pour TRE,
ces mesures ne sont pas définies car un seul état est présent dans l'espace
commun.

\begin{figure}[H]
\centering
\begin{subfigure}[t]{0.48\linewidth}
\centering
\safeincludegraphics[width=\linewidth]
{outputs_separability/3_silhouette_scores.png}
\caption{Score de silhouette.}
\end{subfigure}
\hfill
\begin{subfigure}[t]{0.48\linewidth}
\centering
\safeincludegraphics[width=\linewidth]
{outputs_separability/4_fisher_hr.png}
\caption{Score de Fisher calculé sur HR.}
\end{subfigure}
\caption{Deux lectures complémentaires de la séparabilité des états.}
\label{fig:silhouette_fisher}
\end{figure}

Le graphique de gauche de la figure~\ref{fig:silhouette_fisher} mesure la
cohésion globale des classes dans l'espace standardisé. Le graphique de
droite se concentre uniquement sur la variable HR. Le fait que les deux
indicateurs donnent le même ordre général renforce l'interprétation : la
difficulté dépend fortement du sport et de la qualité de séparation des
états.

\begin{figure}[H]
\centering
\safeincludegraphics[width=0.94\linewidth]
{outputs_separability/5_feature_ranking_heatmap.png}
\caption{Classement des variables selon leur pouvoir séparateur dans chaque
sport.}
\label{fig:feature_ranking}
\end{figure}

La figure~\ref{fig:feature_ranking} est une carte de chaleur. Chaque ligne
correspond à une variable et chaque colonne à un sport. L'intensité d'une case
représente le rang ou le pouvoir séparateur de la variable. Elle montre que la
même caractéristique n'est pas également utile partout : HR et RR dominent
dans RUN, tandis que les variables respiratoires prennent plus de poids dans
d'autres disciplines. Cette observation justifie l'analyse multi-sport et
montre la limite d'une règle physiologique unique appliquée à toutes les
disciplines.

\subsection{Étude d'ablation de SKA}

Une étude d'ablation retire volontairement un élément du pipeline pour mesurer
son influence. Ici, le pool de test reste composé de RUN, CYC et TRE. Les
modèles sont entraînés une première fois sans SKA, puis une seconde fois avec
SKA.

\begin{figure}[H]
\centering
\safeincludegraphics[width=0.94\linewidth]
{outputs_ablation/ablation_ska_effect.png}
\caption{Variation des performances lorsque SKA est ajouté aux données
d'entraînement.}
\label{fig:ablation}
\end{figure}

Dans la figure~\ref{fig:ablation}, une variation positive signifie que l'ajout
de SKA améliore la métrique ; une variation négative signifie qu'elle baisse.
L'ajout de SKA ne produit pas de changement significatif de l'accuracy globale
des modèles. Le recall de la récupération du Random Forest diminue d'environ
4,9 points avec $p=0{,}028$. Ce résultat doit rester prudent car plusieurs
comparaisons sont réalisées. L'enseignement principal est que, dans cette
expérimentation, SKA est surtout difficile à prédire lui-même et ne dégrade
pas systématiquement tous les autres sports.

\subsection{Importance SHAP et comparaison avec l'ANOVA}

L'importance SHAP globale est obtenue en calculant, pour chaque variable, la
moyenne de la valeur absolue de ses contributions. Une grande valeur signifie
que la variable modifie souvent et fortement les prédictions, quel que soit le
sens de cette modification.

\begin{figure}[H]
\centering
\safeincludegraphics[width=0.84\linewidth]
{outputs_explainability_kg/1_global_feature_importance.png}
\caption{Importance SHAP globale des neuf variables du Random Forest
multi-sport.}
\label{fig:global_shap}
\end{figure}

Dans la figure~\ref{fig:global_shap}, la longueur d'une barre représente
l'influence moyenne de la variable. Les trois premières sont
\texttt{max\_BR}, \texttt{std\_BR} et \texttt{mean\_BR}. Elles sont suivies
par \texttt{max\_HR}, \texttt{min\_RR} et \texttt{mean\_HR}. Ce classement
décrit le comportement du modèle dans le corpus utilisé ; il ne constitue pas
un classement physiologique universel.

\begin{figure}[H]
\centering
\safeincludegraphics[width=0.94\linewidth]
{outputs_validation/1_shap_anova_correlation.png}
\caption{Relation entre l'importance SHAP du modèle et le F-score ANOVA de
chaque variable.}
\label{fig:shap_anova}
\end{figure}

Chaque point de la figure~\ref{fig:shap_anova} correspond à une variable.
L'axe horizontal représente sa séparabilité univariée selon l'ANOVA et l'axe
vertical son utilisation par le Random Forest selon SHAP. Si les deux mesures
donnaient le même classement, les points suivraient une tendance croissante.

La corrélation de Pearson vaut $r=-0,047$ avec $p=0,904$, et la corrélation de
Spearman vaut $\rho=0,267$ avec $p=0,488$. Ces valeurs indiquent l'absence
d'association forte. Ce n'est pas une contradiction : l'ANOVA étudie une
variable isolée, alors que SHAP mesure son rôle dans une décision
multivariée, en interaction avec les autres variables.

\subsection{Cohérence sémantique}

Sur le modèle descriptif ajusté à l'ensemble du corpus, la cohérence moyenne
vaut 0,368 pour le repos, 0,612 pour l'effort et 0,374 pour la récupération.
Ces nombres décrivent la proximité entre les variables qui soutiennent la
prédiction et les profils construits pour la classe annoncée.

\begin{figure}[H]
\centering
\safeincludegraphics[width=0.90\linewidth]
{outputs_explainability_kg/4_consistency_and_hitl.png}
\caption{Cohérence sémantique par état et répartition initiale des catégories
HITL.}
\label{fig:consistency_hitl}
\end{figure}

La figure~\ref{fig:consistency_hitl} met en parallèle le score de cohérence et
la décision opérationnelle. L'effort obtient la cohérence moyenne la plus
élevée, ce qui indique que ses variables importantes sont plus souvent
proches du profil appris. Le repos et la récupération sont plus ambigus. Ces
valeurs sont descriptives et ne doivent pas être confondues avec
l'évaluation out-of-fold finale.

\subsection{Validation finale du mécanisme HITL}

L'évaluation finale réunit les prédictions produites dans les 32 plis LOSO.
L'accuracy poolée est de 0,6995.

\begin{table}[H]
\centering
\caption{Fiabilité out-of-fold selon la décision HITL.}
\label{tab:hitl_oof}
\begin{tabular}{lrrrr}
\toprule
\textbf{Décision} &
\textbf{Fenêtres} &
\textbf{Coverage} &
\textbf{Accuracy} &
\textbf{Risk} \\
\midrule
\textsc{accept} & 871   & 12,8\,\% & \textbf{1,000} & 0,000 \\
\textsc{check}  & 2\,348 & 34,4\,\% & 0,971 & 0,029 \\
\textsc{review} & 3\,602 & 52,8\,\% & 0,450 & 0,550 \\
\bottomrule
\end{tabular}
\end{table}

La ligne \textsc{accept} signifie que 871 fenêtres, soit 12,8\,\% du corpus,
respectent les seuils les plus exigeants. Dans cette expérimentation, elles
sont toutes correctes. La ligne \textsc{check} contient 2\,348 fenêtres avec
97,1\,\% d'accuracy. La catégorie \textsc{review} contient plus de la moitié
du corpus et seulement 45,0\,\% de bonnes prédictions. Le mécanisme concentre
donc les erreurs dans le groupe destiné à l'examen humain.

\begin{figure}[H]
\centering
\safeincludegraphics[width=0.82\linewidth]
{outputs_hitl_reliability/1_accuracy_by_hitl.png}
\caption{Accuracy out-of-fold de chacune des trois catégories HITL.}
\label{fig:hitl_accuracy}
\end{figure}

Dans la figure~\ref{fig:hitl_accuracy}, la hauteur des barres diminue de
\textsc{accept} vers \textsc{review}. Cette hiérarchie est le comportement
attendu d'un système sélectif. La taille des groupes doit cependant être lue
avec le tableau : une accuracy élevée sur une petite catégorie ne signifie
pas que tout le corpus est automatiquement fiable.

Parmi les 871 fenêtres \textsc{accept}, 868 proviennent de TRE. L'accuracy de
1,000 ne constitue donc pas une garantie uniforme pour RUN, CYC et SKA. Elle
montre surtout que le score sélectionne un sous-ensemble TRE particulièrement
homogène dans cette version expérimentale.

\begin{figure}[H]
\centering
\safeincludegraphics[width=0.92\linewidth]
{outputs_hitl_reliability/2_accuracy_by_sport_and_hitl.png}
\caption{Répartition et accuracy des catégories HITL selon le sport.}
\label{fig:hitl_by_sport}
\end{figure}

La figure~\ref{fig:hitl_by_sport} décompose les résultats précédents. Elle
montre que la catégorie \textsc{accept} est surtout alimentée par TRE, que RUN
et CYC possèdent principalement des prédictions \textsc{check}, et que SKA est
orienté vers \textsc{review}. Ce graphique est indispensable pour éviter une
lecture trop générale du taux de 100\,\%.

\begin{figure}[H]
\centering
\begin{subfigure}[t]{0.48\textwidth}
\centering
\safeincludegraphics[width=\linewidth]
{outputs_hitl_reliability/3_risk_coverage_curve.png}
\caption{Risk en fonction de la coverage.}
\end{subfigure}
\hfill
\begin{subfigure}[t]{0.48\textwidth}
\centering
\safeincludegraphics[width=\linewidth]
{outputs_hitl_reliability/4_coverage_accuracy_tradeoff.png}
\caption{Accuracy en fonction de la coverage.}
\end{subfigure}
\caption{Compromis entre quantité de prédictions automatisées et fiabilité.}
\label{fig:risk_coverage}
\end{figure}

Les courbes de la figure~\ref{fig:risk_coverage} sont obtenues en classant les
fenêtres par fiabilité puis en conservant une proportion croissante des
meilleures. Au début, seules les prédictions les plus sûres sont gardées :
l'accuracy est élevée et le risk faible. Lorsque la coverage augmente, des cas
plus difficiles sont inclus et l'accuracy diminue.

Les 25\,\% de prédictions les mieux classées atteignent 100\,\% d'accuracy.
À environ 50\,\% de coverage, l'accuracy reste proche de 98\,\%. À 75\,\% de
coverage, elle descend vers 82\,\%. Ces courbes montrent donc que le score de
fiabilité ordonne les prédictions de manière utile pour répartir le travail
entre automatisation et vérification humaine.

% Difficultés rencontrées

\section{Difficultés rencontrées et solutions apportées}
\label{sec:difficultes}

\subsection{Hétérogénéité de SportDB~2.0}

La première difficulté a été de comprendre l'organisation réelle de
SportDB~2.0. Les fichiers ne suivent pas tous exactement la même structure :
les noms des champs MATLAB, les signaux disponibles et les annotations
changent selon le sport et la séance. Certaines disciplines fournissent
directement HR, RR et BR, alors que d'autres contiennent principalement des
segments ECG.

Une phase d'exploration a donc été nécessaire avant de construire le pipeline.
Des scripts ont été utilisés pour parcourir l'arborescence, afficher les
champs disponibles et comparer les annotations de chaque discipline. Cette
étape a conduit à retenir RUN, CYC, TRE et SKA pour l'étude multi-sport, car
leurs signaux pouvaient être ramenés à une représentation commune. Les sports
ECG-only ont été laissés de côté afin de ne pas introduire une reconstruction
des intervalles RR insuffisamment validée.

\subsection{Harmonisation des protocoles et des états}

Les protocoles sportifs n'utilisent pas les mêmes phases. RUN distingue
plusieurs niveaux d'effort, tandis que d'autres disciplines utilisent des
annotations plus générales. Une fusion directe de toutes les étiquettes aurait
produit des classes incohérentes.

Le choix retenu a été de conserver deux niveaux d'analyse. RUN est étudié
séparément avec six états afin de préserver les transitions fines. Pour
l'étude multi-sport, les annotations sont regroupées en trois états communs :
repos, effort et récupération. Cette solution permet la comparaison entre
disciplines, mais elle entraîne une perte de granularité qui est explicitement
prise en compte dans l'interprétation des résultats.

\subsection{Prévention des fuites de données}

Le fenêtrage glissant produit plusieurs observations proches pour un même
sujet. Une séparation aléatoire aurait pu placer des fenêtres très similaires
à la fois dans le train et dans le test, donnant une estimation trop
optimiste de l'accuracy.

Pour éviter cette fuite de données, la validation est réalisée par groupe avec
un protocole LOSO. Le scaler, les profils sémantiques et les explications SHAP
sont recalculés dans chaque pli à partir du train uniquement. Cette contrainte
a rendu l'implémentation plus longue, mais elle était nécessaire pour obtenir
une évaluation plus réaliste sur des sujets non observés.

\subsection{Comparaison de modèles de nature différente}

Le Random Forest et le MLP utilisent des variables statistiques, alors que le
LSTM, le GRU et le Transformer exploitent directement des séquences. La
comparaison n'était donc pas immédiate : les modèles devaient partager les
mêmes groupes de test, le même espace de classes et des conditions
d'entraînement aussi proches que possible.

Les architectures temporelles ont été gardées compactes et l'early stopping a
été utilisé pour limiter le surapprentissage et le temps de calcul. Les mêmes
plis LOSO ont été conservés pour tous les modèles. Les résultats ont ensuite
été comparés avec l'accuracy globale, mais aussi avec le recall par classe,
car un gain local sur un état important peut être masqué par la moyenne.

\subsection{Construction de la cohérence sémantique}

SHAP indique quelles variables soutiennent une prédiction, mais ne dit pas
directement si cette explication est cohérente avec l'état annoncé. Il a donc
fallu définir une mesure supplémentaire sans utiliser d'information provenant
du groupe de test.

La solution retenue consiste à construire, dans chaque pli, des profils
statistiques à partir du train. Les principales contributions SHAP positives
sont ensuite comparées à ces profils. La confiance du modèle et la cohérence
sémantique sont combinées pour produire les catégories \textsc{accept},
\textsc{check} et \textsc{review}. Les seuils restent une preuve de concept :
ils servent à étudier le comportement du système et devront être recalibrés
avant toute utilisation réelle.

\subsection{Interprétation de résultats parfois contradictoires}

Une dernière difficulté a été d'interpréter des résultats qui ne donnent pas
toujours le même classement. Le GRU obtient la meilleure moyenne sur RUN, le
Random Forest est premier en multi-sport, et le Transformer présente un
meilleur recall sur certaines classes. De la même manière, SHAP et l'ANOVA ne
classent pas les variables dans le même ordre.

Ces différences n'ont pas été considérées comme des erreurs. Elles ont conduit
à compléter l'accuracy par des analyses par classe, des tests statistiques et
une étude de séparabilité. Cette démarche a permis de distinguer l'effet de
l'architecture de celui de la qualité des données.


% Discussion

\section{Discussion : comparaison avec les travaux existants}
\label{sec:discussion}

\subsection{Comparaison avec les travaux sur la surveillance sportive}

Claudino et al. montrent que l'intelligence artificielle est principalement
utilisée dans le sport pour la prédiction de performance et du risque
~\cite{claudino2019}. Notre travail se situe dans le même objectif général
d'aide à la décision, mais se concentre sur une tâche différente : reconnaître
l'état fonctionnel présent dans une portion de signaux cardio-respiratoires.
La contribution n'est donc pas un score global de performance ou de blessure,
mais une chaîne qui suit des phases successives d'une séance.

Estrella et Capdevila utilisent la variabilité de la fréquence cardiaque, des
modèles classiques et SHAP pour identifier des profils sportifs
~\cite{estrella2025}. Comme eux, nous constatons que le Random Forest constitue
un compromis intéressant entre performance et explicabilité. Notre approche
ajoute cependant trois dimensions : la comparaison avec des modèles
temporels, la prédiction d'états successifs et une décision HITL fondée sur la
confiance et la cohérence.

\subsection{Comparaison avec les travaux sur la validation entre sujets}

Dehghani et al. insistent sur le risque d'une validation aléatoire lorsque des
fenêtres chevauchantes d'un même sujet se retrouvent dans le train et le test
~\cite{dehghani2019}. Notre protocole suit cette recommandation en séparant les
groupes pendant l'évaluation.

Notre contribution va plus loin que la seule accuracy : le scaler, les
profils sémantiques, les valeurs SHAP et la catégorie HITL sont également
recalculés dans chaque pli. Ainsi, l'explication et la fiabilité sont évaluées
dans les mêmes conditions que la prédiction. Cette extension out-of-fold est
importante, car une explication produite sur une fenêtre déjà vue pourrait
sembler artificiellement plus cohérente.

\subsection{Comparaison des modèles tabulaires et temporels}

Les résultats ne montrent pas de domination globale significative du LSTM, du
GRU ou du Transformer. Sur RUN, le GRU obtient la meilleure accuracy moyenne,
mais la variabilité entre groupes empêche une conclusion générale. Ce résultat
rappelle qu'une architecture plus complexe n'est pas automatiquement meilleure
lorsque le nombre de groupes est limité.

L'analyse par classe apporte néanmoins une différence importante. Les modèles
temporels doublent presque le recall de l'effort intense par rapport aux
modèles tabulaires. Leur intérêt apparaît donc surtout dans les transitions,
où l'ordre des valeurs peut contenir une information perdue par les moyennes
et les écarts-types.

Dans l'étude multi-sport, les quatre phases d'effort de RUN sont fusionnées.
Cette simplification rend la cible commune aux sports, mais supprime une
partie de la granularité sur laquelle les réseaux temporels étaient utiles.
La comparaison RUN/multi-sport constitue ainsi une contribution analytique :
elle montre que le choix des classes peut modifier l'avantage attendu d'une
architecture.

\subsection{Comparaison avec les travaux d'explicabilité en sport}

La revue de Kranzinger et al. montre que SHAP domine actuellement l'XAI en
science du sport, tandis que les explications fondées sur des connaissances
du domaine et les évaluations centrées sur l'utilisateur restent limitées
~\cite{kranzinger2025}. La revue de Shen et al. aboutit à un constat proche
pour les signaux physiologiques : l'explicabilité locale et globale progresse,
mais la validation et la mise en contexte restent difficiles
~\cite{shen2025}.

Notre pipeline reprend SHAP parce qu'il fournit une attribution locale et une
importance globale. Il ne s'arrête cependant pas au classement des variables.
Les contributions sont reliées aux concepts HR, BR et RR, puis comparées à des
profils d'état. Le score de cohérence répond donc à une question que SHAP seul
ne traite pas : les variables qui poussent la décision sont-elles proches de
ce qui a été appris pour la classe prédite ?

Cette contribution reste limitée. Le graphe et les profils n'ont pas encore
été validés par un médecin du sport ou un préparateur physique. Notre travail
répond donc partiellement au manque identifié par Kranzinger et al., mais une
évaluation avec des utilisateurs du domaine reste nécessaire.

\subsection{Comparaison avec les graphes de connaissances}

Hogan et al. présentent les graphes de connaissances comme un moyen de
structurer des faits et des relations~\cite{hogan2020}. Rotmensch et al. les
utilisent dans un contexte de santé pour organiser des concepts issus de
données médicales~\cite{rotmensch2017}. Lecue insiste sur leur rôle potentiel
dans l'explication de systèmes d'intelligence artificielle
~\cite{lecue2020}.

Notre graphe est plus léger qu'une ontologie médicale. Il ne contient ni
raisonnement clinique complet ni règles universelles. Son apport est
opérationnel : il traduit une variable technique comme \texttt{std\_BR} en un
concept respiratoire et relie ce concept à l'état fonctionnel. Il constitue
donc une couche intermédiaire entre le calcul SHAP et l'explication destinée à
l'utilisateur.

\subsection{Comparaison avec la classification sélective et le HITL}

Geifman et El-Yaniv proposent de sélectionner les prédictions à partir de la
confiance du modèle~\cite{geifman2017}. Wu et al. montrent plus généralement
que le HITL peut attribuer à un humain les cas pour lesquels le modèle est
insuffisant~\cite{wu2021}.

Notre approche conserve ce principe, mais ne repose pas uniquement sur la
probabilité maximale. Une prédiction doit également présenter une cohérence
sémantique suffisante. Cette combinaison est l'une des principales
contributions du TER : un modèle très confiant peut être orienté vers
\textsc{review} si les arguments qui soutiennent sa décision sont éloignés du
profil appris.

Les résultats out-of-fold montrent une séparation nette entre
\textsc{accept}, \textsc{check} et \textsc{review}. Ils doivent cependant être
lus avec la distribution par sport. La majorité des \textsc{accept} provient
de TRE ; le système sélectionne donc bien un sous-ensemble facile, mais ne
garantit pas la même couverture pour chaque discipline.

\subsection{Apport global de notre contribution}

Par rapport aux travaux étudiés, notre contribution combine dans un même
prototype :

\begin{itemize}[leftmargin=1.5em]
  \item des signaux physiologiques issus de plusieurs sports ;
  \item une comparaison entre modèles tabulaires et temporels sur les mêmes
  fenêtres de 50 échantillons ;
  \item une validation par groupe ;
  \item une analyse des performances et de la séparabilité ;
  \item une explication locale et globale avec SHAP ;
  \item une représentation sémantique des variables ;
  \item une cohérence calculée à partir des profils du train ;
  \item une décision HITL évaluée out-of-fold.
\end{itemize}

L'originalité se situe donc dans l'intégration de ces éléments, plutôt que
dans la création d'un nouveau classifieur. Cette chaîne permet de passer de
« le modèle prédit effort » à une sortie plus complète : « le modèle prédit
effort avec telle confiance, principalement à cause de telles variables, avec
tel niveau de cohérence, et recommande telle action humaine ».

\subsection{Limites}

Les principales limites sont les suivantes :

\begin{itemize}[leftmargin=1.5em]
  \item le corpus multi-sport ne contient que 32 groupes dans la version
  expérimentale ;
  \item les protocoles, les capteurs et les durées diffèrent entre sports ;
  \item TRE ne contient qu'un seul état dans l'espace commun ;
  \item l'harmonisation à trois états masque la granularité de RUN ;
  \item certaines variables respiratoires sont absentes ou peu informatives
  selon les disciplines ;
  \item le module SHAP explique le Random Forest et non directement le LSTM,
  le GRU ou le Transformer ;
  \item les profils sémantiques sont appris dans le corpus et ne constituent
  pas des normes médicales ;
  \item les seuils HITL sont des choix de preuve de concept ;
  \item les graphiques et explications n'ont pas encore été évalués avec des
  entraîneurs ou des spécialistes du sport ;
  \item aucune validation externe sur un second corpus n'a été réalisée.
\end{itemize}

% Conclusion

\section{Conclusion et perspectives}
\label{sec:conclusion}

\subsection{Bilan}

Ce TER a permis de concevoir une chaîne complète allant des signaux
cardio-respiratoires à une prédiction explicable et contrôlable. Le travail ne
s'est pas limité à comparer cinq modèles. Il a également étudié la
séparabilité des données, les contributions des variables, leur mise en
relation avec des concepts physiologiques et la sélection des cas nécessitant
une vérification humaine.

Sur RUN, aucune architecture ne domine significativement l'accuracy globale.
Les modèles temporels améliorent toutefois nettement le recall de l'effort
intense. Ce résultat montre que l'ordre des 50 échantillons peut apporter une
information réelle pour certaines transitions.

L'étude multi-sport met en évidence une forte influence de la discipline et de
la séparabilité des variables. Le Random Forest obtient la meilleure moyenne
LOSO dans l'expérimentation, mais les écarts entre groupes restent importants.
L'analyse des graphes permet de voir que la difficulté ne provient pas
uniquement du choix du modèle.

SHAP fournit les variables qui soutiennent chaque décision. Le graphe de
connaissances traduit ces variables en concepts HR, BR et RR. La cohérence
sémantique complète ensuite la confiance prédictive. Enfin, le mécanisme HITL
classe les fenêtres en \textsc{accept}, \textsc{check} ou \textsc{review} et
concentre les erreurs dans la catégorie destinée au contrôle humain.

\subsection{Réponse à la problématique}

Il est possible de construire un système qui reconnaît des états
physiologiques dans des données multi-sport tout en affichant les raisons
principales de la décision et son niveau de fiabilité. La qualité de cette
reconnaissance dépend toutefois du protocole, de la définition des états, du
nombre de groupes et de la séparabilité des signaux.

La contribution principale n'est donc pas un classifieur universel. Il s'agit
d'une méthode intégrée pour :

\[
\boxed{
\begin{gathered}
\text{prédire}
\;\longrightarrow\;
\text{expliquer}
\;\longrightarrow\;
\text{vérifier la cohérence}\\[0.15cm]
\longrightarrow\;
\text{orienter la décision humaine}
\end{gathered}
}
\]

Cette chaîne rend la sortie du modèle plus lisible et limite l'automatisation
aux cas considérés comme les plus fiables dans l'expérimentation.

\subsection{Perspectives}

Plusieurs prolongements sont possibles :

\begin{itemize}[leftmargin=1.5em]
  \item intégrer les sports disponibles uniquement sous forme d'ECG avec une
  détection validée des pics R ;
  \item ajouter des signaux de mouvement, de vitesse ou d'accélération ;
  \item construire un espace d'états hiérarchique qui conserve les
  sous-phases propres à chaque sport ;
  \item appliquer des méthodes d'explicabilité directement adaptées aux LSTM,
  GRU et Transformers ;
  \item enrichir le graphe avec des règles validées par un spécialiste du
  sport ;
  \item présenter les explications dans une interface permettant de consulter
  l'état, la confiance, les variables principales et la décision HITL ;
  \item calibrer les seuils selon le coût réel d'une erreur et la quantité de
  travail humain acceptable ;
  \item évaluer la compréhension des graphes et des explications auprès
  d'entraîneurs, d'athlètes ou de professionnels de santé ;
  \item valider le pipeline sur un corpus externe.
\end{itemize}

% Annexe

\appendix

\section{Principaux scripts du projet}

\begin{table}[H]
\centering
\caption{Scripts principaux et rôle dans le pipeline.}
\begin{tabularx}{\linewidth}{lX}
\toprule
\textbf{Script} & \textbf{Rôle} \\
\midrule
\texttt{explore\_dataset.py}
& Exploration de l'arborescence, des champs MATLAB et des annotations. \\

\texttt{deep\_learning\_run.py}
& Comparaison RF, MLP, LSTM, GRU et Transformer sur RUN à six états. \\

\texttt{deep\_learning\_multisport.py}
& Construction du corpus commun et évaluation LOSO multi-sport. \\

\texttt{feature\_distribution\_analysis.py}
& Distributions des variables et projection PCA. \\

\texttt{separability\_analysis.py}
& ANOVA, score de Fisher, silhouette et classement des variables. \\

\texttt{ablation\_sport\_inclusion.py}
& Étude de l'effet de l'inclusion de SKA dans le train. \\

\texttt{explainability\_kg.py}
& SHAP, profils d'état, graphe de connaissances et cohérence sémantique. \\

\texttt{explainability\_kg\_oof.py}
& Explications et décisions HITL recalculées dans chaque pli LOSO. \\

\texttt{hitl\_reliability\_analysis\_oof.py}
& Analyse des catégories HITL et courbes risk--coverage. \\
\bottomrule
\end{tabularx}
\end{table}

% Bibliographie

\newpage
\begin{thebibliography}{99}

\bibitem{linardatos2021}
P.~Linardatos, V.~Papastefanopoulos et S.~Kotsiantis,
``Explainable AI: A Review of Machine Learning Interpretability Methods'',
2021.\\
\url{https://www.mdpi.com/1099-4300/23/1/18}

\bibitem{ribeiro2016}
M.~T.~Ribeiro, S.~Singh et C.~Guestrin,
``Why Should I Trust You?: Explaining the Predictions of Any Classifier'',
2016.\\
\url{https://arxiv.org/abs/1602.04938}

\bibitem{lundberg2017}
S.~M.~Lundberg et S.-I.~Lee,
``A Unified Approach to Interpreting Model Predictions'',
2017.\\
\url{https://arxiv.org/abs/1705.07874}

\bibitem{lundberg2019}
S.~M.~Lundberg, G.~Erion, H.~Chen et al.,
``Explainable AI for Trees: From Local Explanations to Global Understanding'',
2019.\\
\url{https://arxiv.org/abs/1905.04610}

\bibitem{breiman2001}
L.~Breiman,
``Random Forests'',
2001.\\
\url{https://www.stat.berkeley.edu/~breiman/randomforest2001.pdf}

\bibitem{goodfellow2016}
I.~Goodfellow, Y.~Bengio et A.~Courville,
\emph{Deep Learning},
2016.\\
\url{https://www.deeplearningbook.org}

\bibitem{cho2014}
K.~Cho, B.~van Merriënboer, C.~Gulcehre et al.,
``Learning Phrase Representations using RNN Encoder--Decoder for Statistical
Machine Translation'',
2014.\\
\url{https://arxiv.org/abs/1406.1078}

\bibitem{vaswani2017}
A.~Vaswani, N.~Shazeer, N.~Parmar et al.,
``Attention Is All You Need'',
2017.\\
\url{https://arxiv.org/abs/1706.03762}

\bibitem{hogan2020}
A.~Hogan, E.~Blomqvist, M.~Cochez et al.,
``Knowledge Graphs'',
2021.\\
\url{https://arxiv.org/abs/2003.02320}

\bibitem{rotmensch2017}
M.~Rotmensch, Y.~Halpern, A.~Tlimat, S.~Horng et D.~Sontag,
``Learning a Health Knowledge Graph from Electronic Medical Records'',
2017.\\
\url{https://www.nature.com/articles/s41598-017-05778-z}

\bibitem{claudino2019}
J.~G.~Claudino et al.,
``Current Approaches to the Use of Artificial Intelligence for Injury Risk
Assessment and Performance Prediction in Team Sports'',
2019.\\
\url{https://link.springer.com/article/10.1186/s40798-019-0202-3}

\bibitem{sportdb2data}
L.~Burattini, S.~Romagnoli, A.~Sbrollini, M.~Morettini, A.~Nocera,
D.~Bondi, T.~Pietrangelo et V.~Verratti,
\emph{Sport DB 2.0},
2024.\\
\url{https://data.mendeley.com/datasets/kzkjkt7mx2/2}

\bibitem{sportdb2paper}
S.~Romagnoli, A.~Sbrollini, A.~Nocera, M.~Morettini, D.~Bondi,
T.~Pietrangelo, V.~Verratti et L.~Burattini,
``Sport DB 2.0: A New Database of Data Acquired by Wearable and Portable
Devices While Practicing Sport'',
2023.\\
\url{https://cinc.org/archives/2023/pdf/CinC2023-067.pdf}

\bibitem{dehghani2019}
A.~Dehghani, T.~Glatard et E.~Shihab,
``Subject Cross Validation in Human Activity Recognition'',
2019.\\
\url{https://arxiv.org/abs/1904.02666}

\bibitem{demsar2006}
J.~Demšar,
``Statistical Comparisons of Classifiers over Multiple Data Sets'',
2006.\\
\url{https://www.jmlr.org/papers/volume7/demsar06a/demsar06a.pdf}

\bibitem{elgendi2014}
M.~Elgendi,
``Revisiting QRS Detection Methodologies for Portable, Wearable,
Battery-Operated, and Wireless ECG Systems'',
2014.\\
\url{https://journals.plos.org/plosone/article?id=10.1371/journal.pone.0084018}

\bibitem{geifman2017}
Y.~Geifman et R.~El-Yaniv,
``Selective Classification for Deep Neural Networks'',
2017.\\
\url{https://arxiv.org/abs/1705.08500}

\bibitem{kranzinger2025}
S.~Kranzinger, C.~Halmich, D.~Hofer et C.~Kranzinger,
``A Scoping Review of Explainable Artificial Intelligence in Sports Science'',
2025.\\
\url{https://link.springer.com/article/10.1007/s44163-025-00709-8}

\bibitem{shen2025}
J.~Shen, J.~Wu, H.~Liang et al.,
``Physiological Signal Analysis Using Explainable Artificial Intelligence:
A Systematic Review'',
2025.\\
\url{https://doi.org/10.1016/j.neucom.2024.128920}

\bibitem{estrella2025}
T.~Estrella et L.~Capdevila,
``Identification of Athleticism and Sports Profiles Throughout Machine
Learning Applied to Heart Rate Variability'',
2025.\\
\url{https://www.mdpi.com/2075-4663/13/2/30}

\bibitem{lecue2020}
F.~Lecue,
``On the Role of Knowledge Graphs in Explainable AI'',
2020.\\
\url{https://www.semantic-web-journal.net/system/files/swj2198.pdf}

\bibitem{wu2021}
X.~Wu, L.~Xiao, Y.~Sun, J.~Zhang, T.~Ma et L.~He,
``A Survey of Human-in-the-Loop for Machine Learning'',
2021.\\
\url{https://arxiv.org/abs/2108.00941}

\bibitem{shlens2014}
J.~Shlens,
``A Tutorial on Principal Component Analysis'',
2014.\\
\url{https://arxiv.org/abs/1404.1100}

\bibitem{rousseeuw1987}
P.~J.~Rousseeuw,
``Silhouettes: A Graphical Aid to the Interpretation and Validation of Cluster Analysis'',
1987.\\
\url{https://wis.kuleuven.be/stat/robust/papers/publications-1987/rousseeuw-silhouettes-jcam-sciencedirectopenarchiv.pdf}

\bibitem{gu2012}
Q.~Gu, Z.~Li et J.~Han,
``Generalized Fisher Score for Feature Selection'',
2012.\\
\url{https://arxiv.org/abs/1202.3725}

\bibitem{nistAnova}
NIST/SEMATECH,
``One-Way ANOVA Overview'',
2013.\\
\url{https://www.itl.nist.gov/div898/handbook/prc/section4/prc431.htm}

\end{thebibliography}

\end{document}
